\documentclass[11pt]{amsart}
% Draft for Comm. Math. Phys.; converted to the Springer Nature template at submission.
\usepackage[T1]{fontenc}
\usepackage[utf8]{inputenc}
\usepackage{lmodern}
\usepackage{amsmath,amssymb,amsthm,mathtools,thmtools}
\usepackage[margin=1.15in]{geometry}
\usepackage{enumitem}
\usepackage{booktabs}
\usepackage[dvipsnames]{xcolor}
\usepackage[colorlinks=true,linkcolor=MidnightBlue,citecolor=MidnightBlue,urlcolor=MidnightBlue]{hyperref}
\usepackage[capitalise,nameinlink]{cleveref}

% ---- theorem environments (thmtools, so that cleveref names each kind correctly)
% main_pt.tex (the Portuguese translation) defines \versaopt before reading this file.
\ifdefined\versaopt
\declaretheorem[name=Teorema,refname={Teorema,Teoremas},numberwithin=section]{theorem}
\declaretheorem[name=Proposição,refname={Proposição,Proposições},sibling=theorem]{proposition}
\declaretheorem[name=Lema,refname={Lema,Lemas},sibling=theorem]{lemma}
\declaretheorem[name=Corolário,refname={Corolário,Corolários},sibling=theorem]{corollary}
\declaretheorem[name=Definição,refname={Definição,Definições},sibling=theorem,style=definition]{definition}
\declaretheorem[name=Observação,refname={Observação,Observações},sibling=theorem,style=remark]{remark}
\declaretheorem[name=Teorema Principal,numbered=no]{mainthm}
\declaretheorem[name=Teorema,numbered=no]{introthm}
\else
\declaretheorem[name=Theorem,numberwithin=section]{theorem}
\declaretheorem[name=Proposition,sibling=theorem]{proposition}
\declaretheorem[name=Lemma,sibling=theorem]{lemma}
\declaretheorem[name=Corollary,sibling=theorem]{corollary}
\declaretheorem[name=Definition,sibling=theorem,style=definition]{definition}
\declaretheorem[name=Remark,sibling=theorem,style=remark]{remark}
\declaretheorem[name=Main Theorem,numbered=no]{mainthm}
\declaretheorem[name=Theorem,numbered=no]{introthm}
\fi

% ---- drafting aids (remove before submission)
\newcommand{\todo}[1]{{\color{BrickRed}[\textbf{TODO:} #1]}}
\newcommand{\src}[1]{{\color{Gray}\footnotesize[\texttt{#1}]}}

% ---- notation
\newcommand{\R}{\mathbb{R}}
\newcommand{\C}{\mathbb{C}}
\newcommand{\Z}{\mathbb{Z}}
\renewcommand{\Re}{\operatorname{Re}}
\renewcommand{\Im}{\operatorname{Im}}
\newcommand{\gbar}{\bar g}
\newcommand{\phibar}{\bar\phi}
\newcommand{\omegabar}{\bar\omega}
\newcommand{\Mhat}{\widehat M}
\newcommand{\Lop}{L}             % linearized RT family
\newcommand{\Cop}{C}             % linearized constraints
\newcommand{\Kop}{K}             % constraint propagation operator
\newcommand{\Lie}{\mathcal{L}}
\newcommand{\sphys}{s_{\ast}}    % the physical root
\newcommand{\sK}{s_{K}}
\newcommand{\sB}{s_{B}}
\newcommand{\Dom}{\mathcal{D}}
\newcommand{\Yp}{Y_{+}}
\newcommand{\mult}{\operatorname{m}}

% DOI links in the bibliography (amsplain prints the note field)
\newcommand{\doi}[1]{\href{https://doi.org/#1}{\texttt{doi:#1}}}


\title[One unstable mode for Choptuik's spacetime]{Exactly one unstable mode: a computer-assisted spectral
count for Choptuik's critical spacetime under spherically symmetric perturbations}
\author{Edivaldo Amaral Gon\c{c}alves}
\address{Instituto Federal de Mato Grosso (IFMT), Cuiab\'a, Mato Grosso, Brazil}
\email{edivaldo.goncalves@ifmt.edu.br}
\urladdr{https://orcid.org/0009-0007-9117-9833}
\thanks{Instituto Federal de Mato Grosso (IFMT), Cuiab\'a, Mato Grosso, Brazil.\newline
E-mail: \texttt{edivaldo.goncalves@ifmt.edu.br}.\quad
ORCID: \href{https://orcid.org/0009-0007-9117-9833}{0009-0007-9117-9833}.}
\date{Draft of \today}

\begin{document}

\begin{abstract}
Choptuik's spacetime, the discretely self-similar solution at the threshold of black hole formation in the
collapse of a massless scalar field, is expected to have exactly one unstable mode, whose growth rate $\lambda$
determines the critical exponent $\gamma=1/\lambda\approx0.374$. The existence of this spacetime was proved by
Reiterer and Trubowitz with a computer-assisted argument. We prove, also with computer assistance, that in their
realization the solution has exactly one unstable mode, modulo gauge, under spherically symmetric linear
perturbations that are smooth in time and analytic in the radial variable up to and across the light cone. The
mode is real and simple, with $\lambda\in[2.674040,\,2.674115]$, so that
$1/\lambda\in[0.373955,\,0.373966]$. More precisely, per Floquet class the linearized operator has exactly four
roots with positive real part: a gauge mode, which translates the time of the singularity; two roots whose
eigenvectors violate the constraints; and the physical mode. The count uses an operator version of Rouch\'e's
theorem with a finite-dimensional reference, a rank-two deflation of the gauge roots, Newton--Kantorovich
enclosures, and bounds in rational and ball arithmetic. The passage from the geometric problem to the operator
rests on a global completeness result for the gauge and an explicit reconstruction of the perturbations. The
certificates, a verification program and the data are publicly available.
\end{abstract}

\maketitle
\setcounter{tocdepth}{1}
\tableofcontents

\section{Introduction}\label{sec:intro}

\subsection{Critical collapse}

In 1993 Choptuik studied numerically the gravitational collapse of a spherically symmetric massless scalar field
and found, at the threshold between dispersion and black hole formation, a universal solution with three striking
features~\cite{Choptuik1993}: it is approached, in all the one-parameter families of initial data that were evolved, when the
parameter is tuned to the threshold; it is discretely self-similar, reproducing itself on smaller and smaller scales with an echoing period
$\Delta\approx3.44$; and near the threshold the mass of the black hole scales like $(p-p_*)^\gamma$ with a universal
exponent $\gamma\approx0.37$. These \emph{critical phenomena} have since been found in many matter models, and are
reviewed in~\cite{GundlachMartinGarcia2007}. For the spherically symmetric Einstein--scalar field system, Christodoulou proved that
sufficiently small data disperse and sufficiently large data form trapped surfaces~\cite{Christodoulou1986,
Christodoulou1991,Christodoulou1993}.

The standard explanation is dynamical~\cite{KoikeHaraAdachi1995,Gundlach1997}. The critical solution is a fixed
point (in the discretely self-similar case, a periodic orbit) of the evolution in logarithmic time, with a stable
manifold of codimension one, which is the threshold itself, and a \emph{single} unstable mode, growing like
$e^{\lambda T}$. A near-critical solution stays close to the critical one until the unstable mode has grown to a
fixed size, and dimensional analysis then gives $\gamma=1/\lambda$, with a periodic modulation in the discretely
self-similar case. Gundlach computed the perturbations of Choptuik's solution numerically and found exactly one growing mode,
with $\gamma=0.374\pm0.001$~\cite{Gundlach1997}, and Mart\'in-Garc\'ia and Gundlach found numerically that all
nonspherical perturbations decay~\cite{MartinGarciaGundlach1999}. A necessary ingredient of this picture is that the
critical solution has exactly one unstable mode; it is this spectral statement that we prove, for spherically
symmetric perturbations.

\subsection{Existence, and the open spectral question}

The existence of Choptuik's spacetime was proved by Reiterer and Trubowitz~\cite{RT2019}, with a computer-assisted
argument in the tradition of Lanford's proof of the Feigenbaum conjectures~\cite{Lanford1982}. They construct a
real analytic, discretely self-similar solution on a neighbourhood of the past light cone of the singularity, as
the fixed point of a contraction near high-precision numerical data, and determine its echoing constant and a
second invariant to $80$ digits. About perturbations they write: ``We do not study perturbations about the
Choptuik spacetime, that are relevant for many of the (conjectured) properties of the solution. Our paper can be
the starting point for a rigorous investigation of this kind. This would be interesting, because current
numerical results seem to be inconclusive''~\cite[\S1]{RT2019}.

This paper carries out such an investigation for spherically symmetric perturbations.

\subsection{The result}

We work with the solution of~\cite{RT2019}, in their coordinates $(\tau,\xi)$ and variables $\omega$, in which the
self-similarity is the translation $\tau\mapsto\tau+2\pi$. Linear perturbations of Floquet type,
$e^{s\tau}$ times a periodic profile, are governed by a family of operators $L(s)$, the linearization of the
Reiterer--Trubowitz system, together with a linearized constraint map $C(s)$. The exponent $s$ corresponds to the
growth rate $\lambda=2\pi s/K$ in the usual logarithmic time, where $e^{-K}$ is the factor by which the
self-similarity rescales lengths.

\begin{introthm}[informal version of the Main Theorem, \cref{sec:main}]
Under spherically symmetric linear perturbations that are smooth in time and analytic in the radial variable up to
and across the light cone, and modulo gauge, Choptuik's spacetime has exactly one unstable mode. The mode is real and algebraically simple, and its growth rate satisfies
\[
  \lambda\in[2.674040,\ 2.674115],\qquad 1/\lambda\in[0.373955,\ 0.373966].
\]
More precisely, per Floquet class the operator $L$ has exactly four roots in the half plane $\Re s>0$, each simple
and with a real representative: a gauge mode, which is the translation of the time of the singularity; two roots whose eigenvectors
violate the constraints; and the physical unstable mode. On the imaginary axis, $L$ has only the root $s=0$, the
translation of the phase. If the light cone of the singularity is held fixed, the gauge mode is no longer removed by the
smaller group of gauge transformations, and the count becomes two; the additional class corresponds to translating
the time of the singularity, not to a second instability.
\end{introthm}

The value of $1/\lambda$ agrees with the measured and computed critical exponent; this is an external check, since
the relation $\gamma=1/\lambda$ concerns the nonlinear evolution. The theorem is spectral: it locates the roots of
$L$ and identifies the physical ones, but says nothing about the linearized evolution semigroup or the nonlinear
dynamics. It concerns spherically symmetric perturbations only, in the regularity class stated above. These
limitations are discussed in \cref{sec:discussion}.

\subsection{Strategy}

The proof has three parts.

\emph{From the geometric problem to the operator} (\cref{sec:physical}). A physical mode is a solution of the
linearized Einstein--scalar field equations of Floquet type, and only its class modulo gauge transformations is
meaningful. We show that every such mode with $\Re s>0$ can be put in the form of~\cite{RT2019} by a gauge
transformation, and then corresponds to an element of $\ker L(s)\cap\ker C(s)$; that conversely every such element
comes from a physical mode; and that the only gauge modes with $\Re s>0$ are the multiples of one explicit mode at
$s=\mu$, the infinitesimal translation of the singular point. The gauge completeness is global on the domain
of~\cite{RT2019} and uses that the light cone lies inside it. The physical multiplicity of a root is then a matter
of linear algebra on its root space.

\emph{Counting} (\cref{sec:counting}). The roots of $L$ with large real part are excluded by a bound on the
numerical range, computed in a function space in which the background acts pointwise. In the remaining bounded
region, the roots are counted with an operator version of Rouch\'e's theorem due to Gohberg and
Sigal~\cite{GohbergSigal1971}: after a rank-two deflation of the two gauge roots, $L$ is compared along the
boundary with a block-diagonal reference whose zeros are the eigenvalues of a $14016\times14016$ matrix, and these
are counted from a Schur decomposition with a certified residual.

\emph{Identification} (\cref{sec:roots}). Each root is enclosed, and shown to be algebraically simple, by a bordered
Newton--Kantorovich argument. The eigenvectors at two of them are shown to violate the constraints by explicit lower
bounds. For the remaining root, the constraints cannot be checked on an approximate eigenvector; instead, an
identity $KC=ML$ propagates the constraints, and a separate certificate shows that the constraint propagation
operator $K$ is injective there.

All the computations are done in exact rational arithmetic, in ball arithmetic, or as a posteriori bounds on
floating-point computations. They take a few days on ordinary computers.

\subsection{What is new}

The tools are classical; what is new is their combination for this operator, and the results. The paper turns into a
theorem a spectral picture that was known only numerically: it gives a rigorous count of the unstable modes of
Choptuik's spacetime; a rigorous enclosure of the
physical growth rate; the classification of all the roots with positive real part; and the passage from the geometric
problem to the operator, including the completeness of the gauge. The certificates, a verification program that
re-checks them, the data and a reproduction guide are publicly available (\cref{sec:computation-data}).

\subsection{Organization}

\Cref{sec:setting} recalls the formulation of~\cite{RT2019} and defines $L$, $C$ and $K$. \Cref{sec:main} states the
Main Theorem. \Cref{sec:physical} relates physical modes to the operator, \cref{sec:counting} counts the roots and
\cref{sec:roots} identifies them and completes the proof. \Cref{sec:computation} describes the computations, their
verification and the use of AI assistants, and \cref{sec:discussion} discusses the limitations and the nonspherical
problem. The appendices contain the proofs of the structural lemmas, explicit formulas and the list of certificates.

\section{The background and the Reiterer--Trubowitz formulation}\label{sec:setting}

This section fixes notation. \Cref{sec:setting-background,sec:setting-system} recall the parts
of~\cite{RT2019} that we use, with their notation wherever possible. \Cref{sec:setting-linear,sec:setting-sectors,sec:setting-constraints,sec:setting-realizations}
define the objects of the Main Theorem: the linearized family $L(s)$, the constraint map $C(s)$, the
constraint propagation operator $K(s)$, the Floquet classes and the function spaces.
Section numbers in references to~\cite{RT2019} are those of its arXiv version 1203.3766v1, whose source we
used throughout.

\subsection{The Choptuik spacetime}\label{sec:setting-background}

Let
\[
  M=\{(u_-,u_+)\in\R^2:\ u_+<0,\ u_+<u_-<-u_+/81\},\qquad
  \xi=\frac{u_+-u_-}{u_++u_-}.
\]
Reiterer and Trubowitz~\cite[Main result]{RT2019} prove the existence of constants $K,\mu>0$, known to
$80$ digits ($K\approx1.72272620$, $\mu\approx0.16830708$), and of real analytic functions
$\phi,\zeta,Q$ on $M$ with $\mu+Q\xi^2>0$, such that the spherically symmetric metric
\begin{equation}\label{eq:metric}
  \gbar=e^{-2\zeta}\Big(-\frac{2\,(du_-\otimes du_++du_+\otimes du_-)}{\mu^2(u_++u_-)^2}
  +\frac{\xi^2}{(\mu+Q\xi^2)^2}\,g_{S^2}\Big)
\end{equation}
on $M\times S^2$ and the scalar field $\phibar=\phi$ solve
$\operatorname{Ric}_g=2\,d\phi\otimes d\phi$, $\Box_g\phi=0$. The map
$\Theta(u_-,u_+)=e^{-2\pi\mu}(u_-,u_+)$ satisfies
\[
  \Theta^*\gbar=e^{-2K}\gbar,\qquad \phibar\circ\Theta=-\phibar,
\]
and the scalar curvature is not identically zero, so the spacetime is not flat and blows up at the
singular point $(u_-,u_+)=(0,0)$, which is not in $M$. The boundary $u_+=u_-$ is the axis of symmetry, a
removable coordinate singularity. The null hypersurface $u_-=0$ is the past light cone of the singular
point; $M$ contains a neighbourhood of it.

\emph{Coordinates.} We use the coordinates $(\tau,\xi)$ of~\cite{RT2019}, defined by
\begin{equation}\label{eq:tauxi}
  u_\sigma=-(1+\sigma\xi)\,e^{-\mu\tau},\qquad\sigma=\pm.
\end{equation}
In them $\Theta$ is the translation $\tau\mapsto\tau+2\pi$, the axis is $\xi=0$, the light cone $u_-=0$ is
$\xi=1$, and the condition $u_-<-u_+/81$ becomes $\xi<41/40$. Replacing $\xi\ge0$ by the Cartesian
variable $x\in\R^3$ with $|x|=\xi$, the spacetime becomes the open set
\[
  \Mhat=\R_\tau\times B_{41/40}(0)\subset\R\times\R^3,
\]
on which $\gbar$ and $\phibar$ are real analytic, including at the axis $x=0$ and across the light cone
$|x|=1$~\cite{RT2019}. All perturbations below are defined on $\Mhat$.

\emph{Logarithmic time.} Since $\Theta$ rescales lengths by $e^{-K}$ and acts by $\tau\mapsto\tau+2\pi$,
the time $T=K\tau/2\pi$ is the usual logarithmic time of critical collapse, in which $\Theta$ is a
translation by $K$. The scalar field returns to itself under $\Theta^2$, so the echoing period is
$\Delta=2K=3.44545240\ldots$, in agreement with the numerical value $3.445452402(3)$ of
Mart\'in-Garc\'ia and Gundlach~\cite{MartinGarciaGundlach2003}, as already noted in~\cite{RT2019}.

\subsection{Frame variables and the system}\label{sec:setting-system}

Reiterer and Trubowitz write the equations as a first order system, quadratically nonlinear, for the
constant $\mu$ and four functions $\omega=(\omega_1,\omega_2,\omega_3,\omega_4)$ of $(\tau,\xi)$. With the
radial null vector fields
\[
  D^\sigma=\sigma\,\partial_\tau+\mu(1+\sigma\xi)\,\partial_\xi,\qquad\sigma=\pm,
\]
and the reflection $(Pf)(\tau,\xi)=f(\tau,-\xi)$, they are
\begin{equation}\label{eq:omega}
\begin{aligned}
  \omega_1&=2\xi Q+(D^-+D^+)\big(\zeta+\log(\mu+\xi^2Q)\big),&
  \omega_2^\sigma&=D^\sigma\big(\zeta+\log(\mu+\xi^2Q)\big)-\sigma\mu,\\
  \omega_3^\sigma&=D^\sigma\zeta-\sigma\mu,&
  \omega_4^\sigma&=D^\sigma\phi,
\end{aligned}
\end{equation}
with $\omega_i^-=\omega_i$ and $\omega_i^+=-P\omega_i$ for $i=2,3,4$~\cite[\S2]{RT2019}. The
Einstein--scalar field equations are equivalent to the vanishing of ten explicit functions
$\Omega_i^\sigma(\mu,\omega)$, $i\in\{1,2,2\sharp,3,4\}$, each of the form
\[
  \Omega_i^\sigma=\xi\,(D^\sigma+\sigma\mu)(\text{one component of }\omega)
  +\mu\,(\text{linear in }\omega)+\xi\,(\text{quadratic in }\omega).
\]
The reflection symmetry reduces them to the five equations $\Omega^+=0$, posed on the cylinder
$(\R/4\pi\Z)\times(-\xi_*,\xi_*)$, $\xi_*>1$, for functions with the symmetries
\begin{equation}\label{eq:symmetries}
  \omega_1,\omega_2,\omega_3\ \text{$2\pi$-periodic},\qquad \omega_4\ \text{$2\pi$-antiperiodic},\qquad
  P\omega_1=-\omega_1.
\end{equation}

\emph{Series.} Each $\omega_i$ is expanded as a Fourier series in $\tau$ with period $4\pi$ and a
Chebyshev series in $\xi$:
\begin{equation}\label{eq:series}
  v(\tau,\xi)=\sum_{m,n\in\Z}v_{mn}\,e^{im\tau/2}e^{in\theta},\qquad \xi=\cos\theta,\qquad
  v_{m,-n}=v_{mn}.
\end{equation}
The symmetries~\eqref{eq:symmetries} say that $\omega_1,\omega_2,\omega_3$ contain only even $m$,
$\omega_4$ only odd $m$, and $\omega_1$ only odd $n$. In these coordinates $\partial_\tau$ acts by
$im/2$, the reflection by $(-1)^n$, multiplication by $\xi$ by averaging the neighbours $n\pm1$, and
products of functions by convolution of coefficients.

\emph{The selected system and the constraints.} Reiterer and Trubowitz split the five equations
$\Omega^+=0$ into a part with as many equations as unknowns and a remainder, the constraints, by two
selection operators $S$ and $S^\sharp$: $\Omega^+=0$ if and only if $S\Omega^+=0$ and $S^\sharp\Omega^+=0$.
The selected part is
\begin{equation}\label{eq:F}
  F(\mu,\omega):=S\,\Omega^+(\mu,\omega)
  =O_\mu\,\omega+\mu\,S\Gamma_1\omega+S\Xi\,\Gamma_2(\omega,\omega),
\end{equation}
where:
\begin{itemize}[leftmargin=*]
\item $O_\mu=\partial_\tau+\mu D_O$, with $D_O=\operatorname{diag}(\xi\partial_\xi+1,\
  (1+\xi)\partial_\xi+1,\ (1+\xi)\partial_\xi+1,\ (1+\xi)\partial_\xi+1)$, is the free part;
\item $\Gamma_1$ is linear and $\Gamma_2$ is symmetric bilinear; both are built from reflections and
  pointwise products, without derivatives;
\item $\Xi$ is multiplication by $\xi$, and $S$ contains the regularized division by $\xi$,
  $R_\xi f=(f-f|_{\xi=0})/\xi$.
\end{itemize}
The explicit formulas are in~\cite[\S3]{RT2019}. We write $F^\sharp(\mu,\omega):=S^\sharp\Omega^+(\mu,\omega)$
for the constraints.

\emph{The background.} Reiterer and Trubowitz construct a solution
$(\mu,\omegabar)$ of $F=0$, $F^\sharp=0$, as a fixed point of a contraction in a weighted $\ell^1$ space
of coefficients with weights $(65/64)^{|m|}(5/4)^{|n|}$, near explicit numerical data, and from it the
spacetime of \cref{sec:setting-background}~\cite[\S\S4--5]{RT2019}. The solution is normalized by the
condition $\operatorname{Re}(\omegabar_4)_{10}=0$, which fixes the translation invariance in $\tau$. Their final
bounds are
\begin{equation}\label{eq:RTball}
  |\mu-(\mu_A+\mu_B)|\le2^{-277},\qquad \|\omegabar-(\omega_A+\omega_B)\|\le2^{-277},\qquad
  \|\omega_B\|\le2^{-25},
\end{equation}
where $(\mu_A,\omega_A)$ and $(\mu_B,\omega_B)$ are explicit data published with~\cite{RT2019}, and the
norm is their weighted $\ell^1$ norm. Most of our bounds use only $(\mu_A,\omega_A)$ and the consequence
$\|\omegabar-\omega_A\|\le2^{-25}+2^{-277}$.

\subsection{The linearized family}\label{sec:setting-linear}

A perturbation of Floquet type is $\omega=\omegabar+\epsilon\,e^{s\tau}h$, where $h$ is a $4\pi$-periodic
multifield. Here $\mu$ is held fixed, and so are the coordinates $(\tau,\xi)$: perturbations of $\mu$ would
move the null coordinates~\eqref{eq:tauxi}. That this is no loss of generality for physical modes is part
of \cref{sec:physical}. The normalization of the phase is not imposed on $h$.

\begin{definition}[The linearized family]\label{def:L}
For $s\in\C$ and a $4\pi$-periodic multifield $h$ with $Ph_1=-h_1$, set
\[
  L(s)h:=e^{-s\tau}\,D_\omega F(\mu,\omegabar)\big[e^{s\tau}h\big],\qquad
  C(s)h:=e^{-s\tau}\,D_\omega F^\sharp(\mu,\omegabar)\big[e^{s\tau}h\big].
\]
\end{definition}

\begin{lemma}[$L$ and $C$ are affine in $s$]\label{lem:affine}
$L(s)=L(0)+s\,I$ and $C(s)=C_0+s\,C_1$, where
\[
  L(0)=J+B,\qquad J=O_\mu,\qquad B=\mu\,S\Gamma_1+2\,S\Xi\,\Gamma_2(\omegabar,\cdot\,),
\]
and $C_0,C_1$ are given explicitly in \cref{app:formulas}.
\end{lemma}

\begin{proof}
In~\eqref{eq:F}, the only derivative in $\tau$ is the one in $O_\mu$, with coefficient one. The operators
$\Gamma_1$, $\Gamma_2(\omegabar,\cdot)$, $\Xi$ and $S$ act in $\xi$ and by pointwise products, so they
commute with multiplication by $e^{s\tau}$, and $e^{-s\tau}\partial_\tau e^{s\tau}=\partial_\tau+s$. In
$\Omega^+=\Xi H_\mu\omega+\mu\Gamma_1\omega+\Xi\Gamma_2(\omega,\omega)$ the derivative $\partial_\tau$ occurs
only in $H_\mu=H_1\partial_\tau+(\text{terms in }\xi)$, where $H_1$ is a constant matrix of identities and
reflections~\cite[\S\S2--3]{RT2019}. Hence the same computation gives $C(s)=C_0+sC_1$ with
$C_1=S^\sharp\Xi H_1$.
\end{proof}

A root of $L$ is a point $s_0$ at which $L(s_0)$ has a nontrivial kernel in the realizations of
\cref{sec:setting-realizations}. By \cref{lem:affine}, the roots of $L$ are the points $s_0$ with
$-s_0\in\sigma(L(0))$, and a Jordan chain of $L$ at $s_0$, that is $L(s_0)h_0=0$ and
$L(s_0)h_j=-h_{j-1}$, is a Jordan chain of the operator $L(0)$ at the eigenvalue $-s_0$. The
\emph{algebraic multiplicity} $\mult(s_0)$ is the dimension of the generalized eigenspace
$\bigcup_k\ker(L(0)+s_0)^k$. It is finite, and the roots are isolated, by \cref{prop:fredholm}.

\subsection{Floquet classes and the two sectors}\label{sec:setting-sectors}

A multifield $h$ is $4\pi$-periodic, so it contains all Fourier indices $m$; the background contains
only the indices allowed by~\eqref{eq:symmetries}. Accordingly, the space of profiles splits as
$X=X_A\oplus X_B$:
\begin{itemize}[leftmargin=*]
\item \emph{sector A}, with the symmetry of the background: $h_1,h_2,h_3$ with even $m$ and $h_4$ with
  odd $m$;
\item \emph{sector B}, with the opposite symmetry: $h_1,h_2,h_3$ with odd $m$ and $h_4$ with even $m$.
\end{itemize}
Since $\omegabar$ belongs to sector A, both sectors are invariant under $L(s)$ and $C(s)$; we write
$L_A,L_B$ and $C_A,C_B$ for the restrictions. In terms of the half period $\tau\mapsto\tau+2\pi$, sector A
consists of the profiles with the same symmetry as the background, and sector B of those with the opposite
one.

Multiplication by $e^{i\tau/2}$ is the shift $m\mapsto m+1$ of the Fourier index. It preserves the
spaces of \cref{sec:setting-realizations} up to a constant factor, commutes with every operation
in $\xi$ and with multiplication by functions of $\tau$, and conjugates $\partial_\tau$ to
$\partial_\tau+i/2$. Hence
\begin{equation}\label{eq:shift}
  e^{-i\tau/2}\,L(s)\,e^{i\tau/2}=L(s+\tfrac{i}{2}),\qquad
  e^{-i\tau/2}\,C(s)\,e^{i\tau/2}=C(s+\tfrac{i}{2}),
\end{equation}
and the shift exchanges the two sectors.

\begin{lemma}[Reduction to sector A]\label{lem:sectorB}
For every $s\in\C$,
\[
  \mult_L(s)=\mult_{L_A}(s)+\mult_{L_A}(s+\tfrac{i}{2}),
\]
and $L_A$ is $i$-periodic: $\mult_{L_A}(s+i)=\mult_{L_A}(s)$. The conjugation~\eqref{eq:shift} maps the
Jordan chains of $L_B$ at $s$ onto those of $L_A$ at $s+\tfrac i2$, and it maps the chains satisfying the
constraints onto chains satisfying the constraints.
\end{lemma}

\begin{proof}
The splitting $X=X_A\oplus X_B$ reduces $L(s)$, so multiplicities add. The shift by one Fourier index
maps $X_B$ onto $X_A$ and, by~\eqref{eq:shift}, conjugates $L_B(s)$ to $L_A(s+\tfrac i2)$ and $C_B(s)$
to $C_A(s+\tfrac i2)$; conjugation preserves kernels, Jordan chains and their lengths. The shift by two
indices maps $X_A$ to itself and gives the $i$-periodicity. The details for the realizations used below,
where the shift is an isometry up to the factor $\kappa_1$, are in \cref{app:lemmas}.
\end{proof}

Consequently the roots of $L$, counted per Floquet class $s+\tfrac{i}{2}\Z$, correspond to the roots of
$L_A$ counted modulo $i\Z$. All computations are done for $L_A$ in the strip
\[
  \widetilde Q=\{-\tfrac14<\Im s\le\tfrac34\},
\]
whose lower part $\{-\tfrac14<\Im s\le\tfrac14\}$ carries the sector-A roots (the \emph{A band}) and whose upper
half carries the sector-B roots, shifted by $\tfrac i2$ (the \emph{B band}). For instance, the root $\sB$
of the Main Theorem appears as $\sB+\tfrac i2$ in the B band.

\emph{Reality.} The coefficients of $\omegabar$ satisfy the reality condition
$v_{-m,-n}=\overline{v_{mn}}$~\cite{RT2019}; hence $\overline{L_A(\bar s)\,\bar h}=L_A(s)h$ for the
corresponding conjugation of coefficients, and the roots of $L_A$ are symmetric under $s\mapsto\bar s$.
For $L_B$ the same argument, combined with the shift, gives the symmetry $s\mapsto\bar s+i$ in the B band
(\cref{app:lemmas}). These symmetries are what make the four roots of the Main Theorem real.

\subsection{Constraint propagation}\label{sec:setting-constraints}

The equations $\Omega^+=0$ are not independent. Reiterer and Trubowitz show that if $F(\mu,\omega)=0$,
then $\omega^\sharp=F^\sharp(\mu,\omega)$ satisfies the linear homogeneous system
\begin{equation}\label{eq:sharp}
  O_\mu^\sharp\,\omega^\sharp+\mu\,R_\xi\Gamma_1^\sharp\,\omega^\sharp+\Gamma_2^\sharp(\omega,\omega^\sharp)=0,
\end{equation}
where $O_\mu^\sharp=\partial_\tau+\mu\operatorname{diag}(\xi\partial_\xi+1,(1+\xi)\partial_\xi+1)$, and
$\Gamma^\sharp_1,\Gamma^\sharp_2$ are again algebraic~\cite[\S3]{RT2019}. Linearizing this identity at
$(\mu,\omegabar)$, without assuming $F(\mu,\omega)=0$, gives the following lemma, proved in
\cref{app:lemmas}.

\begin{lemma}[Propagation of the constraints]\label{lem:KCML}
Let
\[
  K(s)=J^\sharp+B^\sharp+s\,I,\qquad J^\sharp=O^\sharp_\mu,\qquad
  B^\sharp=\mu\,R_\xi\Gamma_1^\sharp+\Gamma_2^\sharp(\omegabar,\cdot\,).
\]
There is an explicit affine family $M(s)=M_0+sM_1$ (\cref{app:formulas}) such that
\begin{equation}\label{eq:KCML}
  K(s)\,C(s)=M(s)\,L(s)
\end{equation}
as closed operators from the domain of $L$ in $Y(a,b)$ to the domain of $K$ in $Y(a',b')$, for all weights
$1<a'<a$, $1<b'<b$ (\cref{sec:setting-realizations}).
\end{lemma}

In particular, if $L(s)h=0$ and $K(s)$ is injective on its domain in the smaller space, then $C(s)h=0$:
an eigenvector of $L$ satisfies the constraints at every $s$ which is not a root of $K$. The loss of
radius from $(a,b)$ to $(a',b')$ is harmless, because all our certificates for $K$ are formulated in the
smaller space $(129/128,9/8)$.

\subsection{Realizations and the Fredholm property}\label{sec:setting-realizations}

We use two families of coefficient spaces. For weights $a,b>1$:
\begin{itemize}[leftmargin=*]
\item $Y(a,b)$ is the weighted $\ell^1$ space of~\cite{RT2019}, with norm
  $\sum_i\eta_i\sum_{m,n}a^{|m|}b^{|n|}|(h_i)_{mn}|$; the component weights $\eta_i>0$ give equivalent
  norms and only affect constants (\cref{app:lemmas}). We write $\Yp=Y(65/64,5/4)$ for the space of the
  existence proof.
\item $\mathbb T(a,b)$, the \emph{torus space}, is the weighted $\ell^2$ space with \emph{signed} Fourier
  index,
  \[
    \|h\|^2_{\mathbb T(a,b)}=\sum_i\sum_{m,n\in\Z}a^{2m}b^{2n}\,|(h_i)_{mn}|^2 .
  \]
  It is $L^2$ of the torus $\{|w|=a\}\times\{|z|=b\}$ in the variables $w=e^{i\tau/2}$, $z=e^{i\theta}$,
  $\xi=(z+z^{-1})/2$. On it, multiplication by an analytic function, the reflection $P$ and the division by
  $\xi$ act pointwise, which makes several bounds sharp (\cref{sec:counting}).
\end{itemize}
The inclusion $Y(a,b)\subset\mathbb T(a,b)$ has norm at most one.

\begin{proposition}[Fredholm property]\label{prop:fredholm}
In each of the spaces $Y(a,b)$ and $\mathbb T(a,b)$ with $1\le a\le65/64$ and $1<b\le5/4$:
\begin{enumerate}[label=\textup{(\roman*)},leftmargin=*]
\item $J$ is closed on its maximal domain $\Dom$, which coincides with the closure of the trigonometric
  polynomials in the graph norm, and $J+z$ is invertible for real $z\ge0$, with compact inverse;
\item $B$ is bounded, so $L(s)=J+B+s$ is closed on the fixed domain $\Dom$ for all $s$;
\item with $Q=(J+z_0)^{-1}$ for a fixed $z_0>0$, the family $L(s)Q=I+(B+s-z_0)Q$ is an analytic
  family of Fredholm operators of index zero, invertible for real $s$ large.
\end{enumerate}
Consequently the roots of $L$ are isolated and have finite algebraic multiplicity. The same holds for
$K(s)$ in the corresponding sharp spaces.
\end{proposition}

The proof is in \cref{app:lemmas}. The key point of~(i) is that the homogeneous solutions of $J+z$ in each
Fourier mode are multiples of $(1+\xi)^{-1-(z+im/2)/\mu}$ or $\xi^{-1-(z+im/2)/\mu}$, which are singular at
$\xi=-1$ or $\xi=0$, inside the Bernstein ellipse; so $J+z$ is injective on the maximal domain.

The certificates of \cref{sec:counting,sec:roots} are computed in the torus space, while the existence
proof, and with it the physical interpretation, lives in $\Yp$. The following lemma shows that the
choice does not matter in the region where we count.

\begin{lemma}[Independence of the realization]\label{lem:Lreal}
For $|s|\le 3.1$, the kernels and the Jordan chains of $L(s)$ in $\Yp$ and in $\mathbb T(65/64,5/4)$
coincide. Hence, in the disc $|s|\le3.1$, the roots of $L$ and their algebraic multiplicities are the same
in the two realizations.
\end{lemma}

\begin{proof}[Sketch]
Write $H(s)=I+(B+s)Q$ with $Q=J^{-1}$, the same operator in both spaces, and split the coefficients into a
finite box $G=\{|m|<1024,\ n<4096\}$ and its complement $T=I-G$. In both spaces
$\|T(H(s)-I)T\|<1$ for $|s|\le3.1$: the bound is $0.371$ in the torus, from a bound
$\|B\|\le3.9642$ certified through the symbol of $B$ (\cref{sec:counting}), and $0.2666$ in $\Yp$, from a bound on the
background part computed from the data and the published ball of~\cite{RT2019} (\cref{app:recovery}); both use closed forms for $\|QT\|$. If $H(s)x=0$ with $x$ in the
torus space, then $Gx$ is a finite vector, $Tx$ is the unique torus solution of
$THT\,Tx=-THGx$, and the Neumann series solves the same equation in $\Yp$; by uniqueness the two
solutions agree, so $x\in\Yp$. Chains are handled by induction, since the right-hand side
$-h_{j-1}$ is already in the strong space. The converse is the inclusion. Details, and the rational
verification of the constants, are in \cref{app:lemmas}.
\end{proof}

\begin{remark}\label{rem:radius}
The same \emph{radius recovery} argument shows more: an element of the kernel, or of a Jordan chain, in
a space of smaller radius automatically lies in the space of larger radius, for every $s$
(\cref{sec:physical}, \cref{lem:recovery}). It is used in \cref{sec:physical} to show that physical modes,
which are only assumed analytic in some unspecified neighbourhood, produce eigenvectors in $\Yp$.
\end{remark}

\section{Main result}\label{sec:main}

Throughout, $(\gbar,\phibar)$ denotes the solution of the Einstein--scalar field equations
$\operatorname{Ric}_{g}=2\,d\phi\otimes d\phi$, $\Box_{g}\phi=0$ constructed by Reiterer and
Trubowitz~\cite{RT2019} on the spherically symmetric manifold $\Mhat$ of \cref{sec:setting}, with
self-similarity constant $K\approx 1.7227$ and null rescaling constant $\mu\approx 0.16831$ (both known to
$80$ digits). As in~\cite{RT2019}, we call it the \emph{Choptuik spacetime}. The discrete self-similarity
$\Theta$ acts by $\tau\mapsto\tau+2\pi$ in the coordinates $(\tau,\xi)$ of RT, rescales lengths by $e^{-K}$
and reverses the sign of $\phibar$; the full period of the background is therefore $\Theta^{2}$, that is,
$\tau\mapsto\tau+4\pi$.

A perturbation is of \emph{Floquet type with exponent $s\in\C$} if it is $e^{s\tau}$ times a profile which is
$4\pi$-periodic in $\tau$ (\cref{def:physical-mode}). Since $e^{i\tau/2}$ is itself $4\pi$-periodic, the
exponent is defined only modulo $\tfrac{i}{2}\Z$; we call the classes of $\C/\tfrac{i}{2}\Z$ \emph{Floquet
classes}, and count everything per class. In terms of the logarithmic time $T$ in which $\Theta$ is the
translation by $K$, the factor $e^{s\tau}$ is $e^{\lambda T}$ with
\begin{equation}\label{eq:lambda}
  \lambda=\frac{2\pi s}{K}.
\end{equation}

The family $L(s)$, $s\in\C$, is the linearization of the Reiterer--Trubowitz system about the background,
acting on Floquet profiles (\cref{def:L}); $C(s)$ is the corresponding linearized constraint map. Physical
Floquet modes and gauge transformations are defined in \cref{sec:physical}
(\cref{def:physical-mode,def:gauge}).

\begin{mainthm}
Let $(\gbar,\phibar)$ be the Choptuik spacetime of~\cite{RT2019}. Consider spherically symmetric linear
perturbations of Floquet type, with $L$ realized in $\Yp$ or, equivalently in the region considered, in the torus
space (\cref{sec:setting-realizations}); physical modes in the regularity class of \cref{def:physical-mode}, smooth in
$\tau$ and real analytic in $\xi$ up to and across the light cone; and gauge transformations as in
\cref{def:gauge}.
\begin{enumerate}[label=\textup{(\arabic*)},leftmargin=*]
\item \textbf{Roots of $L$.} Modulo $\tfrac{i}{2}\Z$, the roots of $L$ in the half plane $\Re s>0$, counted
  with algebraic multiplicity, are exactly four. Each of them is algebraically simple and has a real
  representative:
  \begin{center}
  \begin{tabular}{@{}lll@{}}
  \toprule
  root & enclosure & eigenvector of $L$ \\
  \midrule
  $s=\mu$ & $\mu=0.168307078963\ldots$ (exact) & $g_*=(1+Z)\omegabar$, pure gauge \\
  $s=\sB$ & $|\sB-0.0414194105|\le 3.7\cdot10^{-5}$ & violates the constraints \\
  $s=\sK$ & $|\sK-0.401024733|\le 1.72\cdot10^{-4}$ & violates the constraints \\
  $s=\sphys$ & $\sphys\in[0.73316949,\ 0.73318983]$ & satisfies the constraints \\
  \bottomrule
  \end{tabular}
  \end{center}
  On the imaginary axis, the only root of $L$ is $s\equiv 0$; it is algebraically simple, with eigenvector
  $\partial_\tau\omegabar$, the infinitesimal translation of the phase.
\item \textbf{Physical modes.} The space of physical Floquet modes with $\Re s>0$, modulo gauge, is
  one-dimensional: it consists of the modes with exponent $\sphys$, and there are no generalized modes
  \textup{(}Jordan chains\textup{)} of length $\ge2$. If the past light cone of the singularity is held fixed
  (\cref{def:conefixed}), the quotient by the smaller group of cone-preserving gauge transformations is
  two-dimensional: the additional class has exponent $\mu$ and is tangent to the one-parameter family of
  exact solutions obtained by translating the singular point, all of which are isometric to
  $(\gbar,\phibar)$. It is therefore not a second instability of the critical solution, but the effect of
  fixing the time of the singularity.
\end{enumerate}
\end{mainthm}

Here $Z=\mu^{-1}\partial_\tau+\xi\partial_\xi$. The eigenvector $g_*$ is a profile; the corresponding Floquet
mode $e^{\mu\tau}g_*$ is the Lie derivative of the background along $X_0=\partial_{u_-}+\partial_{u_+}=e^{\mu\tau}(\mu^{-1}\partial_\tau+\xi\partial_\xi)$,
the generator of translations of the singular point $(u_-,u_+)=(0,0)$ (\cref{lem:cone}).

\begin{remark}[The critical exponent]\label{rem:gamma}
By~\eqref{eq:lambda} and the value of $K$ from~\cite{RT2019}, the unique unstable physical mode grows like
$e^{\lambda T}$ with
\[
  \lambda=\frac{2\pi\sphys}{K}\in[2.674040,\ 2.674115],\qquad
  \frac{1}{\lambda}\in[0.373955,\ 0.373966].
\]
The scaling argument of critical collapse relates a single unstable mode to the mass-scaling exponent by
$\gamma=1/\lambda$~\cite{KoikeHaraAdachi1995,Gundlach1997}. Choptuik measured $\gamma\approx0.37$ in his
simulations~\cite{Choptuik1993}. Gundlach computed the unstable mode by numerical perturbation theory and found
$\lambda_1=-2.674\pm0.009$, hence $\gamma=-1/\lambda_1=0.374\pm0.001$~\cite{Gundlach1997}; the sign is that of his
time variable, in which the growing mode has negative exponent, so his $-\lambda_1$ corresponds to our
$\lambda$. Our interval lies within his error bar. This agreement is an external check, not part of the theorem: the
scaling argument concerns the nonlinear evolution, which we do not study.
\end{remark}

\begin{remark}[What is, and is not, asserted]\label{rem:scope}
\begin{enumerate}[label=(\alph*),leftmargin=*]
\item The theorem is \emph{spectral}: it counts Floquet modes of the linearized equations. It asserts
  nothing about the linearized evolution semigroup, nor about the nonlinear dynamics.
\item Only spherically symmetric perturbations are considered. Nonspherical perturbations are discussed
  in \cref{sec:discussion}.
\item Physical modes are taken in the class of \cref{def:physical-mode}: smooth in $\tau$ and real
  analytic in $\xi$, with a complex neighbourhood uniform in $\tau$. Merely smooth dependence on $\xi$
  across the light cone and at the axis is not covered (\cref{sec:discussion}).
\item Reiterer and Trubowitz prove existence of $(\gbar,\phibar)$ near their numerical data; local
  uniqueness is not formally stated in~\cite{RT2019}. The theorem is about their solution.
\item The roots of $L$ which violate the constraints are not physical: the corresponding solutions of
  the evolution part of the system do not solve the Einstein equations. The root $\mu$ is a coordinate
  effect of the moving light cone: its mode is pure gauge, and it survives only in the quotient by the smaller
  group of cone-preserving gauge transformations, which is the second count in~(2).
\end{enumerate}
\end{remark}

\begin{remark}[Neutral modes]\label{rem:neutral}
Part~(1) on the imaginary axis concerns the roots of $L$ only. At $s=0$ there are two further exact
symmetries of the equations, the global scaling $\delta g=2\gbar$ and the shift $\delta\phi=\mathrm{const}$,
which are physical perturbations with $\delta\omega=0$ and are therefore invisible to $L$. The gauge and
reconstruction lemmas of \cref{sec:physical} are proved for $\Re s>0$ only, so part~(2) has no analogue on
the axis.
\end{remark}

\begin{remark}[Dependence on computation]\label{rem:computation}
The proof is computer assisted in two places: the existence of the background with explicit bounds, which
we take from~\cite{RT2019}, and the bounds of \cref{sec:counting,sec:roots}, which are new. All of the
latter are produced in interval arithmetic (Arb~\cite{Arb}) or exact rational arithmetic, or as a posteriori
bounds on floating-point computations with rigorous rounding control. Certificates, the verification
program and the data are publicly available (\cref{sec:computation}).
\end{remark}

\subsection{Structure of the proof}\label{sec:main-structure}
The proof has three layers.
\begin{enumerate}[label=(\roman*),leftmargin=*]
\item \emph{Geometry} (\cref{sec:physical}). Every physical Floquet mode with $\Re s>0$ can be put, by a
  gauge transformation, into the form used by Reiterer and Trubowitz, and then corresponds to an element of
  $\ker L(s)\cap\ker C(s)$; conversely every such element comes from a physical mode. The only gauge modes
  with $\Re s>0$ are the multiples of $e^{\mu\tau}g_*$ (\cref{lem:gauge}). This reduces
  part~(2) to part~(1) together with the behaviour of the constraints on each root space.
\item \emph{Counting} (\cref{sec:counting}). The number of roots of $L$ in the right half plane is
  computed by an operator version of Rouch\'e's theorem with a finite-dimensional reference, after a
  rank-two (Brauer) deflation of the two gauge roots $0$ and $\mu$.
\item \emph{Identification} (\cref{sec:roots}). Each root is localized and shown to be simple by a
  Newton--Kantorovich argument, and its eigenvector is shown either to violate the constraints (by an
  explicit lower bound), or to be pure gauge, or, for $\sphys$, to satisfy them on the whole root space
  (through the injectivity of the constraint propagation operator $K$).
\end{enumerate}

\section{From physical modes to the operator}\label{sec:physical}

Part~(2) of the Main Theorem is a statement about solutions of the linearized Einstein--scalar field
equations, while part~(1) is about the family $L(s)$. This section proves that for $\Re s>0$ the two are
equivalent (\cref{thm:correspondence}), and reduces the physical multiplicity of a root of $L$ to finite
dimensional linear algebra on its root space (\cref{prop:physmult}). The only computer-assisted inputs are
the fact F1 below and the simplicity of the root $\mu$, proved in \cref{sec:roots}.

\subsection{Physical modes and gauge}\label{sec:physical-defs}

We work on $\Mhat=\R_\tau\times B_{41/40}(0)$ with the Cartesian coordinates $(\tau,x)$, $|x|=\xi$, and in
the double null coordinates $u_\pm$ of~\eqref{eq:tauxi} away from the axis. A spherically symmetric
symmetric 2-tensor which is smooth on $\Mhat$ has the form
\begin{equation}\label{eq:invariant-tensor}
  \alpha\,d\tau^2+2\beta\,d\tau\,(x\cdot dx)+\gamma\,(x\cdot dx)^2+d\,|dx|^2,
\end{equation}
with coefficients that are functions of $(\tau,\xi^2)$.

\emph{The RT form.} The background~\eqref{eq:metric} has vanishing components $\gbar_{++}=\gbar_{--}$ in
the null coordinates. A perturbation $(\delta g,\delta\phi)$ is \emph{in RT form} if
$\delta g_{++}=\delta g_{--}=0$; then its two-dimensional part is conformal to $\gbar_{2D}$, and the
linearization of~\eqref{eq:metric} and~\eqref{eq:omega} at fixed $\mu$ and fixed coordinates assigns to it
\emph{RT variables} $\delta\omega=D\Omega(\gbar,\phibar)[\delta g,\delta\phi]$, where
$\Omega:(g,\phi)\mapsto\omega$ is the map~\eqref{eq:omega}. Explicitly, $\delta g=(\delta a/a)\gbar_{2D}+\delta
b\,g_{S^2}$ gives $\delta\zeta=-\delta a/2a$, $\delta W=-W(\delta\zeta+\delta b/2b)$ with $W=\mu+\xi^2Q$,
$\delta Q=\delta W/\xi^2$, and $\delta\omega$ is obtained by differentiating~\eqref{eq:omega}
(\cref{app:lemmas}).

\begin{definition}[Physical Floquet mode]\label{def:physical-mode}
Let $s\in\C$. A \emph{physical Floquet mode with exponent $s$} is a spherically symmetric pair
$\delta=(\delta g,\delta\phi)$ on $\Mhat$ such that:
\begin{enumerate}[label=\textup{(\roman*)},leftmargin=*]
\item $\delta$ solves the Einstein--scalar field equations linearized at $(\gbar,\phibar)$;
\item $(\Theta^2)^*\delta g=e^{4\pi s}e^{-4K}\delta g$ and $(\Theta^2)^*\delta\phi=e^{4\pi s}\delta\phi$;
\item \emph{(regularity)} the coefficients $\alpha,\beta,\gamma,d$ of $\delta g$ in~\eqref{eq:invariant-tensor}
  and $\delta\phi$, multiplied by $e^{-s\tau}$ and by the background factor $e^{2\zeta}$ where it applies
  (so that they are $4\pi$-periodic by~(ii)), are smooth on $\Mhat$ and belong to $C^\infty(\R/4\pi\Z;\mathcal A(E))$
  for some complex neighbourhood $E$ of $[-1,1]$ in the $\xi$-plane, where $\mathcal A(E)$ denotes the bounded
  holomorphic functions on $E$.
\end{enumerate}
A \emph{generalized physical mode of order $k$} is a solution of~(i) of the form
$e^{s\tau}\sum_{j<k}\tau^j\delta_j$, with profiles $\delta_j$ as in~(iii).
\end{definition}

Condition~(iii) asks for smoothness in $\tau$ and real analyticity in $\xi$ up to and across the light cone
$\xi=1$, with a complex neighbourhood that is uniform in $\tau$ but may be arbitrarily thin. The factor
$e^{-4K}$ in~(ii) is the one of the background, $(\Theta^2)^*\gbar=e^{-4K}\gbar$; equivalently,
$\delta\zeta$, $\delta W/W$ and $\delta\phi$ are $e^{s\tau}$ times $4\pi$-periodic functions.

\begin{definition}[Gauge]\label{def:gauge}
A \emph{gauge transformation} is $\delta\mapsto\delta+\Lie_X(\gbar,\phibar)$, where
$\Lie_X(\gbar,\phibar)=(\Lie_X\gbar,X(\phibar))$ and $X$ is a complex vector field on $\Mhat$ which is
either
\begin{enumerate}[label=\textup{(\alph*)},leftmargin=*]
\item of class $C^1$ on $\Mhat$, or
\item continuous on $\Mhat$, of class $C^1$ off the light cone, and such that $\Lie_X(\gbar,\phibar)$ is in
  RT form with analytic RT variables.
\end{enumerate}
The parameter $\mu$ and the coordinates $(\tau,x)$ are held fixed. A physical mode is \emph{pure gauge}
if it equals $\Lie_X(\gbar,\phibar)$ for such an $X$.
\end{definition}

\begin{definition}[Fixed light cone]\label{def:conefixed}
In the \emph{cone-fixed setting}, physical modes are required to leave the light cone null to first order,
$\delta g_{++}=0$ on $\{u_-=0\}$, and gauge transformations are required to be tangent to it,
$X(u_-)=0$ on $\{u_-=0\}$.
\end{definition}

Holding $\mu$ fixed is essential: changing $\mu$ changes the null directions
$\ker(\sigma\,d\xi-\mu(1+\sigma\xi)\,d\tau)$, and a perturbation in RT form for one value of $\mu$ is not in
RT form for another. Fixing $\mu$ and the coordinates is part of the identification of the background.

For $X$ such that $\Lie_X(\gbar,\phibar)$ is in RT form we write
$\delta\omega_X:=D\Omega(\gbar,\phibar)[\Lie_X(\gbar,\phibar)]$. The basic example is
\begin{equation}\label{eq:X0}
  X_0=\partial_{u_-}+\partial_{u_+}=e^{\mu\tau}\big(\mu^{-1}\partial_\tau+\xi\partial_\xi\big),
\end{equation}
the generator of the joint translation of the null coordinates, which moves the singular point $(0,0)$ and
the light cone. It is analytic on $\Mhat$, $\Lie_{X_0}(\gbar,\phibar)$ is in RT form, and
\begin{equation}\label{eq:gaugemode}
  G_*:=\delta\omega_{X_0}=e^{\mu\tau}g_*,\qquad g_*=(1+Z)\,\omegabar,\qquad Z=\mu^{-1}\partial_\tau+\xi\partial_\xi .
\end{equation}
The identity~\eqref{eq:gaugemode} is a direct computation from~\eqref{eq:omega} (\cref{app:lemmas}). Since
$\Lie_{X_0}(\gbar,\phibar)$ solves the linearized equations, step (R4) of \cref{lem:R} below shows that
$e^{\mu\tau}g_*$ is a Floquet mode of $L$ with exponent $\mu$: $L(\mu)g_*=0$ and $C(\mu)g_*=0$.

\subsection{Gauge modes}\label{sec:physical-gauge}

The following fact is checked in exact rational arithmetic from the data of~\cite{RT2019}.

\begin{lemma}[F1]\label{lem:F1}
$\omegabar_4(\cdot,0)\not\equiv0$ and $\omegabar_4(\cdot,-1)\not\equiv0$. More precisely, the Fourier
coefficients of index $m=1$ satisfy $|c_1(\omegabar_4(\cdot,0))|\ge0.493046-3\cdot10^{-8}$ and
$|c_1(\omegabar_4(\cdot,-1))|\ge0.097752-3\cdot10^{-8}$.
\end{lemma}

\begin{proof}
With $\omega(\tau,\xi)=\sum v_{mn}e^{im\tau/2+in\theta}$, $\xi=\cos\theta$, the coefficient is
$c_1=\sum_nv_{1n}i^n$ at $\xi=0$ and $\sum_nv_{1n}(-1)^n$ at $\xi=-1$. For the data $\omega_A$ these sums are
approximately $0.0729160+0.4930463\,i$ and $-0.0794729+0.0977527\,i$, computed exactly. By~\eqref{eq:RTball},
and since the weights of the norm are $\ge1$, $|c_1(\omegabar)-c_1(\omega_A)|\le2^{-25}+2^{-277}<3\cdot10^{-8}$.
\end{proof}

\begin{lemma}[Completeness of gauge]\label{lem:gauge}
Let $X$ be a vector field as in \cref{def:gauge}\textup{(a)} such that $\Lie_X(\gbar,\phibar)$ is spherically
symmetric and in RT form,
and suppose that $\delta\omega_X\neq0$ is of Floquet type with exponent $s$, $\Re s>0$, that is,
$\delta\omega_X\circ\Theta^2=e^{4\pi s}\delta\omega_X$. Then $e^{4\pi(s-\mu)}=1$ and
$\Lie_X(\gbar,\phibar)=c\,\Lie_{X_0}(\gbar,\phibar)$ for some $c\in\C\setminus\{0\}$; in particular
$\delta\omega_X=c\,G_*$. Consequently, the $SO(3)$-average of $X$ is $cX_0$, and $X-cX_0$ is a Killing field of $\gbar$ that
annihilates $\phibar$, for instance a rotation. The same holds for $X$ as in
\cref{def:gauge}\textup{(b)}.
\end{lemma}

In words: the only gauge mode with $\Re s>0$ is $e^{\mu\tau}g_*$, with exponent $\mu$, in sector A.

\begin{proof}
\emph{Step 0 (symmetry).} Averaging $X$ over $SO(3)$ with respect to Haar measure does not change
$\Lie_X(\gbar,\phibar)$, because the background and $\Lie_X(\gbar,\phibar)$ are invariant. The average $\bar X$ is
$C^1$ and invariant, hence of the form $X^\tau\partial_\tau+X^\xi\partial_\xi$ and tangent to the axis, and
$\Lie_{\bar X}(\gbar,\phibar)=\Lie_X(\gbar,\phibar)$. We replace $X$ by $\bar X$; the conclusions below are about
$\bar X$, and $X-\bar X$ has vanishing Lie derivative of the background. (The original field need not be
invariant: rotational Killing fields can be added without changing the perturbation.)

\emph{Step 1 (transport).} In null coordinates, $\gbar_{\pm\pm}=0$ and $\gbar_{+-}\neq0$, so
$(\Lie_X\gbar)_{++}=2\gbar_{+-}\partial_{u_+}X^{u_-}$ and $(\Lie_X\gbar)_{--}=2\gbar_{+-}\partial_{u_-}X^{u_+}$.
The RT form forces $\partial_{u_+}X^{u_-}=\partial_{u_-}X^{u_+}=0$ on $M$. The level sets
$\{u_-=v\}\cap M$ and $\{u_+=w\}\cap M$ are intervals, hence connected, so $X^{u_-}=f_-(u_-)$ and
$X^{u_+}=f_+(u_+)$ on $M$. Since $u_-$ takes all real values on $M$, including $0$ on the light cone,
$f_-\in C^1(\R)$; and $f_+\in C^1((-\infty,0))$.

\emph{Step 2 (axis).} $X$ is tangent to the axis $u_+=u_-$, so $f_+=f_-$ on $(-\infty,0)$. Write $f:=f_-$.

\emph{Step 3 (equivariance).} $\Theta$ preserves the RT form and the frame of~\cite{RT2019} has coefficients
independent of $\tau$, so $\Omega(\Theta^*(g,\phi))=\Omega(g,\phi)\circ\Theta$; and $\Omega$ is invariant under
constant conformal rescalings of $g$. Differentiating, and using $(\Theta^2)^*\gbar=e^{-4K}\gbar$,
$\phibar\circ\Theta^2=\phibar$, we get $\delta\omega_X\circ\Theta^2=\delta\omega_{(\Theta^2)^*X}$. Since
$u_\pm\circ\Theta^2=\Lambda u_\pm$ with $\Lambda=e^{-4\pi\mu}$, the field $(\Theta^2)^*X$ is of the same type,
with $f$ replaced by $\Lambda^{-1}f(\Lambda\,\cdot)$.

\emph{Step 4 (injectivity).} Suppose $\delta\omega_X=0$. Then $D^\sigma(X(\phibar))=\delta\omega_4^\sigma=0$
for $\sigma=\pm$, so $X(\phibar)=c_2$ is constant on the connected set $M$. Since $D^+u_+=D^-u_-=0$ and
$D^+u_-=-D^-u_+=2\mu e^{-\mu\tau}$,
\[
  X(\phibar)=\frac{e^{\mu\tau}}{2\mu}\big(f(u_-)\,\omegabar_4^+-f(u_+)\,\omegabar_4^-\big).
\]
On the axis $u_\pm=-e^{-\mu\tau}$ and $\omegabar_4^\pm(\tau,0)=\mp\omegabar_4(\tau,0)$, which gives
$-\mu^{-1}e^{\mu\tau}f(-e^{-\mu\tau})\,\omegabar_4(\tau,0)=c_2$ for all $\tau$. As $\omegabar_4(\cdot,0)$ is
antiperiodic, it vanishes somewhere, so $c_2=0$. By \cref{lem:F1} it is not identically zero, and being
real analytic it has isolated zeros; hence $f=0$ on a dense subset of $(-\infty,0)$, and by continuity on
$(-\infty,0]$. Outside the cone, $\{\xi>1\}\cap M$, we have $u_+<0$, so $f(u_-)\,\omegabar_4^+=0$ there. If
$f(v)\neq0$ for some $v>0$, then $\omegabar_4^+$ vanishes on a nonempty open set, hence identically by
analyticity, and in particular on the axis, contradicting \cref{lem:F1}. Thus $f\equiv0$ and $X=0$. (For
complex $X$ apply this to the real and imaginary parts.)

\emph{Step 5 (homogeneity).} If $\delta\omega_X\circ\Theta^2=e^{4\pi s}\delta\omega_X$, Steps 3 and 4 give
$(\Theta^2)^*X=e^{4\pi s}X$, that is,
\[
  f(\Lambda u)=\rho\,f(u)\quad(u\in\R),\qquad \rho=e^{4\pi(s-\mu)},\quad |\rho|=e^{4\pi(\Re s-\mu)}.
\]

\emph{Step 6 (regularity at the cone).} $f$ is $C^1$ at $u=0$, because the cone lies inside $M$. Iterating,
$f(\Lambda^ju)=\rho^jf(u)$ with $\Lambda^ju\to0$.
\begin{itemize}[leftmargin=*]
\item If $|\rho|>1$: the left side tends to $f(0)$, so $f(u)=0$.
\item If $|\rho|=1$ and $\rho\neq1$: $\rho^j$ does not converge, so $f(u)=0$. If $\rho=1$: $f(u)=f(0)$.
\item If $\Lambda<|\rho|<1$, i.e.\ $0<\Re s<\mu$: $f(0)=0$, and
  $f'(0)=\lim_jf(\Lambda^ju)/(\Lambda^ju)=\lim_j(\rho/\Lambda)^jf(u)/u$ with $|\rho/\Lambda|>1$, so $f(u)=0$.
\end{itemize}
Since $\delta\omega_X\neq0$, $f\not\equiv0$, so $\rho=1$ and $f$ is constant, i.e.\ $X=cX_0$.

\emph{Variant (b).} Steps 1--5 use only continuity of $f$ at $u=0$, and Step 6 for $\Re s\ge\mu$ as well.
For $0<\Re s<\mu$: $\delta\phi=X(\phibar)$ is continuous with analytic derivatives
$D^\sigma\delta\phi=\delta\omega_4^\sigma$, hence analytic. On the axis,
$-\mu^{-1}e^{\mu\tau}f(-e^{-\mu\tau})\omegabar_4(\tau,0)$ is analytic, so $f$ is analytic on $(-\infty,0)$.
Near a point $(\tau_0,1)$ of the cone with $\omegabar_4^+(\tau_0,1)\neq0$, which exists by \cref{lem:F1},
$f(u_-)=\big(2\mu e^{-\mu\tau}\delta\phi+f(u_+)\omegabar_4^-\big)/\omegabar_4^+$ is analytic, so $f$ is analytic
at $0$. Then $f(\Lambda u)=\rho f(u)$ forces $a_k(\Lambda^k-\rho)=0$ for the Taylor coefficients, so
$f=a_ku^k$ with $s\equiv\mu(1-k)\pmod{\tfrac i2\Z}$, and $\Re s>0$ forces $k=0$.
\end{proof}

\begin{remark}[The role of the light cone]\label{rem:cone-role}
The proof uses that the light cone lies inside the domain and that $X$ is continuous across it. If $X$ were
only required to be regular in $\{\xi<1\}$, then $f(u)=(-u)^{1-s/\mu}$ would produce a ``gauge mode'' for
every $s$. Requiring a cone-preserving gauge means $f(0)=0$; this excludes constant $f$, and then there is
no gauge mode at all with $\Re s>0$.
\end{remark}

The gauge mode $G_*$ has a direct geometric meaning.

\begin{lemma}[The cone]\label{lem:cone}
Let $\psi_\epsilon(u_-,u_+)=(u_-+\epsilon,u_++\epsilon)$, the flow of $X_0$, extended trivially to the
spheres. For small $\epsilon$, $(\gbar_\epsilon,\phibar_\epsilon):=\psi_\epsilon^*(\gbar,\phibar)$ is an exact
solution of the Einstein--scalar field equations on compact subsets of $\Mhat$, isometric to
$(\gbar,\phibar)$, discretely self-similar about the displaced singular point $(-\epsilon,-\epsilon)$, and in RT
form with the same $\mu$ and coordinates. Its RT variables solve all ten equations $\Omega_i^\sigma=0$, and
\[
  \Omega(\gbar_\epsilon,\phibar_\epsilon)=\omegabar+\epsilon\,e^{\mu\tau}g_*+O(\epsilon^2)
\]
uniformly on compact subsets.
\end{lemma}

\begin{proof}
The equations are diffeomorphism invariant. The map $\psi_\epsilon$ preserves $du_\pm$ and the axis
$\{u_+=u_-\}$, so the two-dimensional part of $\gbar_\epsilon$ is still conformal to $du_-du_+$, and
$\gbar_\epsilon$ has the form~\eqref{eq:metric} with $e^{-2\zeta_\epsilon}=e^{-2\zeta\circ\psi_\epsilon}
(u_++u_-)^2/(u_++u_-+2\epsilon)^2$ and $Q_\epsilon$ determined by the area radius; regularity at the axis
follows from the regularity of the metric. The derivative at $\epsilon=0$ is
$D\Omega[\Lie_{X_0}(\gbar,\phibar)]=G_*$ by~\eqref{eq:gaugemode}.
\end{proof}

Thus $e^{\mu\tau}g_*$ is tangent to a curve of exact solutions, all isometric to the background: the same
critical solution with the time $T_*$ of the singularity translated. With the gauge transformations of
\cref{def:gauge} it is pure gauge. In the cone-fixed setting of \cref{def:conefixed} it is not, since $X_0(u_-)=1$, and it
counts as a physical mode; this is the second count in part~(2) of the Main Theorem.

\subsection{Reconstruction: from physical modes to $\ker L\cap\ker C$}\label{sec:physical-forward}

We need a radius recovery statement for all $s$, not only $|s|\le3.1$ as in \cref{lem:Lreal}.

\begin{lemma}[Radius recovery]\label{lem:recovery}
Let $1\le\kappa_1'\le65/64$, $1<\kappa_2'\le9/8$ and $S>0$. For $|s|\le S$, every element of the kernel of
$L(s)$, and every Jordan chain of $L$ at $s$, in the domain of $L$ in $Y(\kappa_1',\kappa_2')$ lies in the
domain of $L$ in $\Yp$.
\end{lemma}

\begin{proof}[Sketch]
As in \cref{lem:Lreal}: with $Q=J^{-1}$, $H(s)=I+(B+s)Q$ and a finite box $G$, it suffices that
$\|T(H(s)-I)T\|<1$ in both spaces. In the weak space this holds because $\|QT\|$ tends to zero with $G$,
uniformly in the Fourier weight (since $Q$ preserves $m$), and $\|B\|$ is controlled by the bound
$\beta_+$ of \cref{app:recovery}, computed from the data and the published ball of~\cite{RT2019}, reweighted from $5/4$
to $\kappa_2'$; in the strong space the same holds with $\beta_+$ itself. Fourier weight $\kappa_1'=1$ is allowed, because the convolution by the background has
norm at most $\sum_k|b_k|\le\sum_k\kappa_1^{|k|}|b_k|$. Then the kernel equation is solved in the strong space
by a Neumann series on $T$, and uniqueness in the weak space identifies the two solutions; chains follow by
induction. Explicit boxes, e.g.\ $G=2^{17}\times2^{20}$ for $S=200$ and $\kappa_2'=1+2^{-6}$, and the
rational verification of the constants are in \cref{app:lemmas}.
\end{proof}

\begin{lemma}[Reconstruction, forward direction]\label{lem:R}
Let $\delta$ be a physical Floquet mode with exponent $s$, $\Re s>0$. Then:
\begin{enumerate}[label=\textup{(\arabic*)},leftmargin=*]
\item there is a gauge transformation, by a vector field $X$ of the same regularity as $\delta$ (in
  particular $C^1$ on $\Mhat$) and of Floquet type, such that $\delta'=\delta+\Lie_X(\gbar,\phibar)$ is in RT
  form;
\item the RT variables of $\delta'$ are $\delta\omega'=e^{s\tau}h$ with $h$ in the domain of $L$ in $\Yp$,
  $L(s)h=0$ and $C(s)h=0$;
\item $h$ is determined by the gauge class of $\delta$ up to the gauge direction
  $\mathcal G_s$, where $\mathcal G_s=\C\,e^{(\mu-s)\tau}g_*$ if $s\equiv\mu\pmod{\tfrac i2\Z}$ and
  $\mathcal G_s=0$ otherwise; and $h\in\mathcal G_s$ if and only if $\delta$ is pure gauge;
\item a generalized physical mode of order $k$ yields in the same way a Jordan chain of $L$ at $s$ of length
  $k$ modulo $\mathcal G_s$, which satisfies the jet constraints~\eqref{eq:jet} below.
\end{enumerate}
\end{lemma}

\begin{proof}
\emph{(R1) Gauge fixing away from resonance.} We look for $X$ with $(\delta g+\Lie_X\gbar)_{\pm\pm}=0$, that
is $\partial_{u_+}X^{u_-}=-\delta g_{++}/(2\gbar_{+-})$ and $\partial_{u_-}X^{u_+}=-\delta g_{--}/(2\gbar_{+-})$.
With $X^{u_\mp}=e^{(s-\mu)\tau}y_\mp$ and $f_\mp:=-\mu e^{-s\tau}\delta g_{\pm\pm}/\gbar_{+-}$, which are
$4\pi$-periodic by \cref{def:physical-mode}(ii) (both $\delta g_{\pm\pm}$ and $\gbar_{+-}$ scale by
$e^{8\pi\mu}e^{-4K}$ under $\Theta^2$, the former with the extra factor $e^{4\pi s}$), the equations become
\[
  T_{\pm1}(s)\,y_\mp=f_\mp,\qquad T_c(s)=\partial_\tau+\mu(\xi-c)\partial_\xi+s-\mu .
\]
For $\Re s>0$ and $s\notin\mu-\tfrac i2\Z$, the unique periodic solution, analytic in $\xi$ near $[-1,1]$,
is
\begin{equation}\label{eq:transport}
  y(\tau,\xi)=(\partial_\tau+s-\mu)^{-1}f(\cdot,c)\,(\tau)
  +\int_0^\infty e^{-(s-\mu)t}\,\big(f-f(\cdot,c)\big)\big(\tau-t,\ c+e^{-\mu t}(\xi-c)\big)\,dt,
\end{equation}
where $(\partial_\tau+s-\mu)^{-1}$ is the Fourier multiplier $1/(s-\mu+im/2)$. The integral converges also
for $0<\Re s<\mu$, because the integrand vanishes at $\xi=c$ and is therefore $O(e^{-\Re s\,t})$; the
characteristic path stays in a convex neighbourhood of $[-1,1]$. Uniqueness follows by expanding a
homogeneous solution in powers of $\xi-c$: the coefficient of degree $n$ solves
$(\partial_\tau+s+\mu(n-1))y_n=0$, which forces $y_n=0$. The formula commutes with $\partial_\tau$, so $y$
has the regularity of $f$. At the axis, the Cartesian regularity of $\delta g$ makes $f_-$ and $Pf_+$ the
same analytic function, and uniqueness gives $y_+=Py_-$, which means that $X$ is regular on the axis.

\emph{(R2) Resonance.} Let $s_0=\mu-im_0/2$ with $\Re s_0=\mu>0$. At $\xi=1$ the equation for $y_-$ reads
$(\partial_\tau-im_0/2)\,y_-(\cdot,1)=f_-(\cdot,1)$, so a periodic solution exists if and only if the
coefficient $\kappa=(f_-(\cdot,1))_{m_0}$ vanishes; when it does, the solution is unique up to adding
$e^{im_0\tau/2}$, i.e.\ up to adding a multiple of $X_0$. Suppose $\kappa\neq0$. Then
$y_-=y_{\mathrm{reg}}+\kappa\tau e^{im_0\tau/2}$, and the same constant appears in $y_+$ because
$f_+(\cdot,-1)=f_-(\cdot,1)$; regularity at the axis requires the same constant in both components. The
gauge field is $X=X_{\mathrm{reg}}+\kappa\tau X_0$, and
\[
  \delta'=A+\tau B,\qquad B=\kappa\Lie_{X_0}(\gbar,\phibar),\qquad
  A=\delta+\Lie_{X_{\mathrm{reg}}}(\gbar,\phibar)+2\kappa\,d\tau\odot X_0^\flat,
\]
both in RT form and of Floquet type. Since $\delta\omega$ is of first order with $D^\sigma\tau=\sigma$, the RT
variables are $\delta\omega'=e^{s_0\tau}(h_1+\tau h_0)$ with $h_0=\kappa e^{im_0\tau/2}g_*\neq0$, and by
\cref{lem:affine} the linearized equations give $L(s_0)h_0=0$ and $L(s_0)h_1=-h_0$. Multiplying by
$e^{-im_0\tau/2}$ (\cref{eq:shift}) and projecting onto sector A, where $g_*$ lies, gives a Jordan chain
of $L_A$ of length two at $s=\mu$, in the domain in $\Yp$ by \cref{lem:recovery}. This contradicts the
algebraic simplicity of the root $\mu$ (\cref{prop:mu}). Hence $\kappa=0$, and the gauge fixing exists and is
unique up to multiples of $X_0$.

\emph{(R3) RT variables.} For $\delta'$ in RT form, the coefficients~\eqref{eq:invariant-tensor} give
$\delta b=d\,\xi^2$, $\delta\zeta=-\tfrac12\mu^2e^{2\zeta}(\gamma\xi^2+d)$ and
$\delta Q=-\tfrac12We^{2\zeta}\big(d(2\mu Q+\xi^2Q^2)-\mu^2\gamma\big)$, which are even and analytic in $\xi$.
Hence $\delta\omega'=e^{s\tau}h$ with $h$ of class $C^\infty(\R/4\pi\Z;\mathcal A(E'))$ for a slightly smaller
neighbourhood $E'$; the encoding $\delta\omega_i^+=-P\delta\omega_i^-$ is automatic.

\emph{(R4) Equations.} The definitions~\eqref{eq:omega} imply, after linearization, the identities
$\Omega_1^\sigma=\Omega_2^\sigma+\Omega_{2\sharp}^\sigma$ and $\Omega_j^-=\Omega_j^+$ for $j=2,3,4$, and the
curvature formulas of~\cite[\S2]{RT2019} express the linearized Einstein--scalar field equations through the
remaining combinations of $\Omega$'s. Since $\delta'$ solves the linearized equations, all $\delta\Omega_i^\sigma$
vanish off the axis, and by continuity on the axis. In particular $L(s)h=0$ and $C(s)h=0$.

\emph{(R5) The spectral space.} The Fourier coefficients of $h$ decay faster than any power of $|m|$, and
its Chebyshev coefficients decay like $\varrho^{-n}$ for the parameter $\varrho>1$ of a Bernstein ellipse
inside $E'$. Hence $h\in Y(1,\kappa_2')$ for $1<\kappa_2'<\min(\varrho,9/8)$, and $Jh=-(B+s)h$ lies in the same
space, so $h$ is in the domain of $L$ there. By \cref{lem:recovery}, $h$ is in the domain in $\Yp$; in particular $e^{s\tau}h$ extends analytically to
$|\xi|<41/40$. Beyond the light cone, for $1<\xi<41/40$, both families of characteristics $D^\pm$ satisfy
$d\xi/d\tau>0$, so the hypersurfaces $\{\xi=\mathrm{const}\}$ are spacelike, and a solution of the linearized
system there is determined by its Cauchy data on $\{\xi=1+\epsilon\}$, uniformly in the $\tau$-circle of the
periodic profile. Since $\delta\omega'$ and $e^{s\tau}h$ agree near $[-1,1]$ and are both smooth solutions, they
agree on all of $\Mhat$. The same applies to $\delta'$ itself in (R6).

\emph{(R6) Uniqueness modulo gauge.} If $\delta\omega'=0$, then~\eqref{eq:omega} give $\delta Q=0$ and
$\delta\zeta,\delta\phi$ constant; a constant is not of Floquet type with $\Re s>0$, so $\delta'=0$ and $\delta$
is pure gauge. If $\delta\omega'=c\,G_*$ (in the label $s$), then $\delta+\Lie_{X-cX_0}(\gbar,\phibar)$ has
$\delta\omega=0$ and the same conclusion holds. Two gauge fixings differ by an admissible field, whose
$\delta\omega$ is in $\mathcal G_s$ by \cref{lem:gauge}.

\emph{(R7) Chains.} For a generalized mode, the transport is solved order by order,
$T_c(s)y_j=f_j-y_{j-1}$, with gauge fields polynomial in $\tau$. Away from resonance every order is solvable.
At resonance the compatibility condition at each order is forced as in (R2): an extra secular term would
lengthen a Jordan chain of $L$ at $\mu$, which is simple. Steps (R3)--(R5) apply to each profile and give
\[
  \delta\omega'=e^{s\tau}\sum_{j=0}^{k-1}\frac{\tau^j}{j!}\,h_{k-1-j}.
\]
Since $\partial_\tau$ enters the selected equations and the constraints only through the constant coefficients
$1$ and $C_1$ (\cref{lem:affine}), the linearized selected operator maps $e^{s\tau}\tau^jh$ to
$e^{s\tau}\big(\tau^jL(s)h+j\tau^{j-1}h\big)$, and the linearized constraint map maps it to
$e^{s\tau}\big(\tau^jC(s)h+j\tau^{j-1}C_1h\big)$. Collecting powers of $\tau$, the linearized equations become
$L(s)h_i=-h_{i-1}$ and the constraints $C(s)h_i+C_1h_{i-1}=0$ ($h_{-1}=0$): $(h_0,\dots,h_{k-1})$ is a Jordan chain of
$L$ at $s$, in the domain in $\Yp$, satisfying the jet constraints~\eqref{eq:jet}. For $k=2$ this is the form
$e^{s\tau}(h_1+\tau h_0)$ of (R2). If
$h_0\in\mathcal G_s$, then by (R6) the lowest order is pure gauge, and the chain shortens modulo gauge.
\end{proof}

\subsection{The converse}\label{sec:physical-converse}

\begin{lemma}[Reconstruction, converse direction]\label{lem:V}
Let $\Re s>0$, and let $h$ be a $4\pi$-periodic multifield with $Ph_1=-h_1$, in the domain of $L$ in
$\Yp$, with $L(s)h=0$ and $C(s)h=0$. Then there is a unique $\delta=(\delta g,\delta\phi)$ in RT form whose RT
variables are $e^{s\tau}h$, and $\delta$ is a physical Floquet mode with exponent $s$, real analytic on
$\Mhat$ including the axis and the light cone. The map $h\mapsto\delta$ is linear and injective.
\end{lemma}

\begin{proof}
\emph{All ten equations.} The selected components of $\delta\Omega^+$ vanish at $\xi=0$ for every $h$ with $h_1$
odd, because their principal and quadratic parts carry a factor $\xi$; so nothing is lost in the division
$R_\xi$ contained in $S$. Hence $L(s)h=0$ gives $\tfrac12(1+P)\delta\Omega_1^+=\delta\Omega_2^+=\delta\Omega_3^+=
\delta\Omega_4^+=0$, and $C(s)h=0$ gives $\tfrac12(1-P)\delta\Omega_1^+=\delta\Omega_{2\sharp}^+=0$. Thus
$\delta\Omega^+=0$, and $P\delta\Omega_i^\sigma=-\delta\Omega_i^{-\sigma}$ gives $\delta\Omega^-=0$.

\emph{Potentials.} We look for $(\delta Q,\delta\zeta,\delta\phi)=e^{s\tau}(q,z,p)$ with
$\delta\omega=e^{s\tau}h$. By~\eqref{eq:omega} and~\eqref{eq:metric} this means
$D^\sigma_s z=h_3^\sigma$, $D_s^\sigma p=h_4^\sigma$ and $D_s^\sigma(z+\delta\log W)=h_2^\sigma$, with
$D_s^\sigma=D^\sigma+\sigma s$, together with the algebraic relation for $q$ from $\omega_1$. With
$a_i=\tfrac12[(1-\xi)h_i^+-(1+\xi)h_i^-]$ and $b_i=(h_i^++h_i^-)/(2\mu)$, the equation $D_s^\sigma z=h_i^\sigma$ for
both signs is equivalent to $(\partial_\tau+s)z=a_i$ and $\partial_\xi z=b_i$. A direct computation gives
\[
  -2\mu\xi\big((\partial_\tau+s)b_i-\partial_\xi a_i\big)=\delta\Omega_j^--\delta\Omega_j^+,\qquad (i,j)=(3,3),(4,4),(2,2),
\]
so the two equations are compatible. For $\Re s>0$, $(\partial_\tau+s)$ is invertible on $4\pi$-periodic
functions, with inverse $\int_0^\infty e^{-st}a(\tau-t,\xi)\,dt$, and $z=(\partial_\tau+s)^{-1}a_i$ then also
satisfies $\partial_\xi z=b_i$. There is no period obstruction, because $e^{4\pi s}\neq1$. The constraint
$C(s)h=0$ is needed for the identity behind $\log W$; the other compatibilities follow from $L(s)h=0$.

\emph{Geometry.} With $(q,z,p)$ so defined, $\delta$ is in RT form, has RT variables $e^{s\tau}h$, satisfies
the Floquet condition by construction, and solves the linearized equations by the curvature formulas of
\cref{lem:R}~(R4), since all $\delta\Omega$'s vanish. As $h\in\Yp$, the profiles are analytic in $\tau$ in a
strip and in $\xi$ inside the Bernstein ellipse of parameter $5/4$, whose real semi-axis is
$\tfrac12(\tfrac54+\tfrac45)=\tfrac{41}{40}$; $q,z,p$ are even in $\xi$, so $\delta$ is real analytic on $\Mhat$ in
Cartesian coordinates. Injectivity is clear since $\delta$ determines its RT variables.
\end{proof}

\subsection{Correspondence and physical multiplicity}\label{sec:physical-count}

Combining \cref{lem:gauge,lem:R,lem:V}:

\begin{theorem}[Correspondence]\label{thm:correspondence}
Let $\Re s>0$. The map $\delta\mapsto h$ of \cref{lem:R} induces a linear isomorphism
\[
  \frac{\{\text{physical Floquet modes with exponent }s\}}{\{\text{pure gauge modes}\}}
  \;\xrightarrow{\ \cong\ }\;
  \frac{\ker L(s)\cap\ker C(s)}{\mathcal G_s},
\]
with $\mathcal G_s$ as in \cref{lem:R}. Generalized physical modes of order $k$ correspond to Jordan chains of
$L$ at $s$ of length $k$ satisfying the jet constraints
\begin{equation}\label{eq:jet}
  C(s)h_0=0,\qquad C(s)h_j+C_1h_{j-1}=0\quad(1\le j<k),
\end{equation}
modulo $\mathcal G_s$.
\end{theorem}

Note that the jet constraints are not the conditions $C(s)h_j=0$ for each $j$; imposing the latter can
give the wrong dimension already for $2\times2$ examples.

The constraints are handled on the whole root space at once. Fix a root $s_0$ with $\Re s_0>0$, let $E$ be its
root space (the generalized eigenspace of $L(0)$ at $-s_0$, finite dimensional and contained in the domain in
$\Yp$), and $A=L(0)|_E$. Comparing powers of $s$ in~\eqref{eq:KCML}, with $K(s)=K(0)+s$, gives $M_1=C_1$,
$M_0=C_0+K(0)C_1-C_1L(0)$ and $K(0)C_0=M_0L(0)$. Define
\[
  \mathsf D=C_0-C_1L(0)\quad\text{on }E .
\]
Then $K(0)\mathsf D=\mathsf DA$, so $\ker(\mathsf D|_E)$ is $A$-invariant, and for a Jordan chain
$Ah_j=-s_0h_j-h_{j-1}$ one has $\mathsf Dh_j=C(s_0)h_j+C_1h_{j-1}$. Thus a chain satisfies the jet
constraints~\eqref{eq:jet} if and only if it lies in $\ker(\mathsf D|_E)$.

\begin{proposition}[Physical multiplicity]\label{prop:physmult}
Let $s_0$ be a root of $L$ with $\Re s_0>0$ and root space $E$. The dimension of the space of physical
generalized modes with exponent $s_0$, modulo pure gauge, is
\[
  \dim E-\operatorname{rank}(\mathsf D|_E)-\dim(\mathcal G_{s_0}\cap E).
\]
In particular, if $s_0$ is algebraically simple with eigenvector $h$, this number is $1$ if $C(s_0)h=0$ and
$h\notin\mathcal G_{s_0}$, and $0$ otherwise. In the cone-fixed setting of \cref{def:conefixed}, the same holds without the last
term.
\end{proposition}

\begin{proof}
By \cref{thm:correspondence}, the physical generalized modes modulo gauge correspond to
$\ker(\mathsf D|_E)/(\mathcal G_{s_0}\cap E)$, and rank--nullity gives the formula; here $\mathcal G_{s_0}\cap E\subset\ker(\mathsf D|_E)$,
because $\mathsf Dg_*=C(\mu)g_*=0$. The quotient does not
require $\mathcal G_{s_0}$ to have an $A$-invariant complement. For a simple root, $E=\C h$ and
$\mathsf Dh=C(s_0)h$.

In the cone-fixed setting, let $\delta$ be a physical mode with $\delta g_{++}=0$ on the cone. Then
$f_-(\cdot,1)=0$ in (R1), so $y_-(\cdot,1)$ solves $(\partial_\tau+s-\mu)y_-(\cdot,1)=0$; it vanishes away from
resonance, and at resonance the free multiple of $e^{im_0\tau/2}$ is set to zero. The gauge fixing is
therefore cone-preserving and unique, and the obstruction $\kappa$ of (R2) vanishes automatically. There
are no cone-preserving gauge modes with $\Re s>0$: in the proof of \cref{lem:gauge} the condition
$f(0)=0$ excludes the constant solution (\cref{rem:cone-role}). Conversely, the modes of \cref{lem:V} are in
RT form, so $\delta g_{++}=0$ everywhere. The correspondence of \cref{thm:correspondence} thus holds with
$\mathcal G_s$ replaced by $0$; for $s=\mu$ the element $g_*$ corresponds to $\Lie_{X_0}(\gbar,\phibar)$,
which is a physical mode but not a cone-preserving gauge mode.
\end{proof}

\Cref{prop:physmult} reduces part~(2) of the Main Theorem to part~(1) and to three facts about the
eigenvectors: those at $\sB$ and $\sK$ violate the constraints, the one at $\mu$ spans $\mathcal G_\mu$, and the
one at $\sphys$ satisfies them. These are proved in \cref{sec:roots}.

\section{Counting the roots}\label{sec:counting}

This section proves the counting part of the Main Theorem. All statements are for the torus realization
$\mathbb T(65/64,5/4)$ of $L_A$ and, for $K$, the sharp torus $\mathbb T^\sharp(129/128,9/8)$; by
\cref{lem:Lreal} the counts are the same in the realization $\Yp$ of~\cite{RT2019}. We write
\[
  \Omega_A=(0,3)\times(-\tfrac14,\tfrac14),\qquad \Omega_B=(0,3)\times(\tfrac14,\tfrac34)
\]
for the A and B bands of \cref{sec:setting-sectors}, truncated at $\Re s=3$, and $N(L,\Omega)$ for the
number of roots in $\Omega$ counted with algebraic multiplicity.

\begin{theorem}[Counts]\label{thm:counts}
\begin{enumerate}[label=\textup{(\roman*)},leftmargin=*]
\item $L(s)$ is invertible for $\Re s>2.9261$ and $K(s)$ is invertible for $\Re s>1.765$, for all $\Im s$.
\item $N(L_A,\Omega_A)=3$ and $N(L_A,\Omega_B)=1$.
\item $L_A$ has no roots on $\partial\Omega_A\cup\partial\Omega_B$ other than $s=0$, which is an algebraically
  simple root.
\end{enumerate}
\end{theorem}

\begin{corollary}\label{cor:four}
Modulo $\tfrac i2\Z$, $L$ has exactly four roots in $\{\Re s>0\}$, counted with algebraic multiplicity. On the
imaginary axis its only root is $s\equiv0$, which is algebraically simple.
\end{corollary}

\begin{proof}
By \cref{lem:sectorB}, the roots of $L$ modulo $\tfrac i2\Z$ correspond, with multiplicities, to the roots of
$L_A$ modulo $i\Z$, which are counted in the strip $\widetilde Q=\{-\tfrac14<\Im s\le\tfrac34\}$. By~(i)
there are none with $\Re s\ge3$, and by~(iii) none on the lines $\Im s=\tfrac14$ and $\Im s=-\tfrac14\equiv\tfrac34$
with $0<\Re s<3$. So the roots with $\Re s>0$ are those in $\Omega_A\cup\Omega_B$, and there are $3+1$ of
them. On the axis, the segment $\{\Re s=0,\ -\tfrac14\le\Im s\le\tfrac34\}$ is a full period of $L_A$ and lies
in $\partial\Omega_A\cup\partial\Omega_B$, so~(iii) applies.
\end{proof}

The proof of \cref{thm:counts}~(i) is in \cref{sec:counting-large}. Parts (ii) and (iii) are proved in
\cref{sec:counting-rouche,sec:counting-reference,sec:counting-windows,sec:counting-finite} by an
operator version of Rouch\'e's theorem: after a rank-two deflation of the roots $0$ and $\mu$, $L_A$ is
compared on $\partial\Omega$ with a reference whose zeros are the eigenvalues of a $14016\times14016$
matrix, and these are counted from a Schur decomposition with certified residual.

\subsection{Large real part}\label{sec:counting-large}

On the torus space the background operators are pointwise. With $w=e^{i\tau/2}$ and $z=e^{i\theta}$ on the
torus $|w|=a$, $|z|=b$, multiplication by a field $f$ becomes multiplication by its analytic extension
$f(w,z)$, the reflection $P$ becomes $z\mapsto-z$, and the division by $\xi$ becomes multiplication by
$1/\xi(z)$, $\xi(z)=(z+z^{-1})/2$, which is bounded since $b>1$. Consequently
$\Re\langle -Bx,x\rangle$ is the integral over the torus of a Hermitian form in the seven values
$(x_1(z),x_2(\pm z),x_3(\pm z),x_4(\pm z))$ (using $x_1(-z)=-x_1(z)$), with coefficients given by the background,
and
\[
  \max\Re W(-B)\le\sup_{\text{torus}}\lambda_{\max}\big(N^{-1/2}\operatorname{Herm}(-\mathcal S)N^{-1/2}\big),
\]
where $\mathcal S$ is the $7\times7$ symbol and $N$ accounts for the multiplicity of the values. Here
$W(T)=\{\langle Tx,x\rangle:\|x\|=1\}$ is the numerical range. This bound captures all cancellations
between components and Fourier phases.

\begin{lemma}[Numerical range of the background]\label{lem:F3}
In $\mathbb T(65/64,5/4)$, $\max\Re W(-B)\le2.9492$ and $\|B\|\le3.9642$; for the free part,
$\max\Re W(-J)\le-0.02319$. In $\mathbb T^\sharp(129/128,9/8)$, $\max\Re W(-B^\sharp)\le1.4396$ and
$\|B^\sharp\|\le1.6497$, and $\max\Re W(-K(0))\le 1.765$.
\end{lemma}

\begin{proof}[Method of proof]
The supremum of $\lambda_{\max}$ of the symbol over the torus is bounded on a grid of $2048\times4096$
centres for $L$ and $1024\times2048$ for $K$. At each centre the symbol is evaluated in ball arithmetic (Arb), and the inequality is verified by a
$LDL^{\mathsf T}$ factorization in balls of the shifted matrix. Between centres, Lipschitz constants of the
background fields, computed in exact rational arithmetic, give the slack, also evaluated in ball arithmetic.
The Lipschitz slack added between centres is at most $0.0772$ for $L$ and $0.0536$ for $K$. The free part is handled by a closed
form. Details and file references are in \cref{sec:computation}.
\end{proof}

\begin{proof}[Proof of \cref{thm:counts}~(i)]
For $x$ in the domain, $\Re\langle L(s)x,x\rangle\ge(\Re s-2.9492+0.02319)\|x\|^2$, so
$\|L(s)x\|\ge(\Re s-2.9261)\|x\|$, and $L(s)$ is injective for $\Re s>2.9261$. Since $L(s)$ is Fredholm of
index zero (\cref{prop:fredholm}), it is invertible. The argument for $K$ is the same.
\end{proof}

The bound $2.9261$ should be compared with the largest root, $\sphys\approx0.733$: the exclusion is not sharp,
which is why the rectangles $\Omega_A,\Omega_B$ extend to $\Re s=3$. Before the symbol was used, block-by-block
majorants gave a bound larger by two orders of magnitude.

\subsection{Rouch\'e's theorem for operator families, and deflation}\label{sec:counting-rouche}

We use the operator version of Rouch\'e's theorem of Gohberg and Sigal~\cite{GohbergSigal1971}, in its
homotopy form. Let $\Omega$ be a bounded domain with piecewise smooth boundary, and $H,R$ analytic families of
Fredholm operators of index zero on a neighbourhood of $\overline\Omega$. Suppose moreover that $H-R$ is compact. If
$R+t(H-R)$ is invertible on $\partial\Omega$ for every $t\in[0,1]$, then $H$ and $R$ have the same number of zeros in $\Omega$, counted with
algebraic multiplicity. A sufficient condition is $\|R(s)^{-1}(H(s)-R(s))\|<1$ on $\partial\Omega$.

\emph{Normalization.} The free part $J+s$ is triangular in the Chebyshev index within each Fourier mode, with
diagonal $\mu(n+1)+s+im/2$, so $Q_0(s):=(J+s)^{-1}$ is analytic and invertible in $\Re s>-\mu$. The family
$H(s)=\widetilde L(s)Q_0(s)$ has the same zeros as $\widetilde L$, with the same multiplicities, and
$H(s)=I+(B+E)Q_0(s)$, with $E=\widetilde L-L$ the deflation below, is the identity plus a compact operator.

\emph{Deflation.} The roots $0$ and $\mu$ are known exactly: $L(s)x_0=s\,x_0$ for $x_0=\partial_\tau\omegabar$
(translation invariance in $\tau$), and $L(s)x_1=(s-\mu)\,x_1$ for $x_1=g_*$ (\cref{sec:physical-defs}). The
root $0$ lies on $\partial\Omega_A$, and the root $\mu$ keeps the resolvent large between $0$ and $\mu$. We
remove both by a finite-rank perturbation.

\begin{lemma}[Brauer deflation]\label{lem:brauer}
Let $\phi_0,\phi_1$ be bounded functionals and $\delta_0,\delta_1\in\C$, and set
$\widetilde L(s)=L(s)+\delta_0x_0\phi_0^*+\delta_1x_1\phi_1^*$. Then
\[
  \det\big(I+L(s)^{-1}(\widetilde L(s)-L(s))\big)=\frac{q(s)}{s(s-\mu)},
\]
where $q$ is the quadratic polynomial $\det\big(\operatorname{diag}(s,s-\mu)+[\delta_j\phi_i^*x_j]_{ij}\big)$, and
for every $s_0$ the algebraic multiplicities satisfy
\[
  \mult_{\widetilde L}(s_0)=\mult_L(s_0)+\operatorname{ord}_{s_0}\frac{q(s)}{s(s-\mu)} .
\]
\end{lemma}

\begin{proof}
Since $L(s)^{-1}x_0=x_0/s$ and $L(s)^{-1}x_1=x_1/(s-\mu)$, the perturbation $L(s)^{-1}(\widetilde L-L)$ has
rank two with range in $\operatorname{span}\{x_0,x_1\}$, and the determinant reduces to a $2\times2$
determinant. The multiplicity formula is the logarithmic-residue form of the Gohberg--Sigal theory for the
factorization $\widetilde L=L\,(I+L^{-1}(\widetilde L-L))$, in which the second factor is a finite-rank
meromorphic perturbation of the identity whose logarithmic multiplicity is the order of its determinant.
\end{proof}

We take $\phi_i$ supported in the finite box $Z$ below, biorthogonal to explicit approximations $x_j^A$ computed
from the data $\omega_A$ ($\phi_i^*x_j^A=\delta_{ij}$), and $(\delta_0,\delta_1)=(1,1+\mu_A)$. With
$e_{ij}=\phi_i^*(x_j-x_j^A)$,
\[
  q(s)=(s+1+e_{00})\big(s-\mu+\delta_1(1+e_{11})\big)-\delta_0\delta_1e_{01}e_{10},
\]
so the two deflated roots move to approximately $-1$. The errors satisfy $\|P_Z(x_0-x_0^A)\|\le1.5\cdot10^{-9}$
and $\|x_1-x_1^A\|\le2.51\cdot10^{-8}$, so $|e_{ij}|\le\|\phi_i\|\,\|x_j-x_j^A\|\le5.51\cdot2.51\cdot10^{-8}$, and both factors of $q$ have modulus
$\ge1-10^{-6}$ in $\Re s\ge0$. Hence $q$ has no zeros in $\Re s\ge0$, and by \cref{lem:brauer}
\begin{equation}\label{eq:brauer-count}
  N(L_A,\Omega_A)=N(\widetilde L_A,\Omega_A)+1,\qquad N(L_A,\Omega_B)=N(\widetilde L_A,\Omega_B),
\end{equation}
the $+1$ coming from the pole of $q(s)/(s(s-\mu))$ at $\mu\in\Omega_A$. Moreover, at every $s$ with $\Re s\ge0$ where
$\widetilde L_A(s)$ is invertible, $\mult_L(s)=1$ if $s\in\{0,\mu\}$ and $\mult_L(s)=0$ otherwise. In the certificates the
deflation uses $x_j^A$, and the differences $x_j-x_j^A$ are carried as perturbations.

\subsection{The reference and its verification on the contour}\label{sec:counting-reference}

Partition the coefficient space of sector A into
\begin{itemize}[leftmargin=*]
\item the finite box $Z=\{|m|\le31,\ 0\le n<128\}$, of dimension $14016$;
\item $26$ \emph{windows} $J$ of at most $16$ consecutive Fourier modes, covering the rest of the box
  $\{|m|<200,\ n<400\}$;
\item the \emph{far} part, the complement.
\end{itemize}
Let $P_Z,P_J,P_{\mathrm{far}}$ be the coordinate projections and $\widehat H_{ab}(s)=P_aH(s)P_b$. The reference is the
block-diagonal family
\[
  R(s)=\operatorname{diag}\big(\widehat H_{ZZ}(s),\ \widehat H_{JJ}(s)\ (J\text{ windows}),\ I_{\mathrm{far}}\big).
\]
It is analytic in $\Re s>-\mu$, and its zeros in $\Omega$ are those of the blocks. In the certificates the finite
block is taken one step further from the operator: with $W$ the diagonal matrix of the weights, $\widehat H_{ZZ}$ is
replaced by $R_Z(s)=W^{-1}(\mathsf A+s)W\,(P_Z(J+s)P_Z)^{-1}$, where $\mathsf A$ is the stored floating-point matrix that
approximates $W\widetilde L_Z(0)W^{-1}$, built from the data $\omega_A$ and the vectors $x_j^A$ and taken as an exact
matrix. The difference $P_Z\widetilde LP_Z-W^{-1}\mathsf AW$, due to the background error, the deflation vectors and
the rounding in the assembly, is part of $H-R$ and is bounded in the entries below.

\begin{lemma}[The finite block]\label{lem:ZZ}
$\widehat H_{ZZ}(s)=\widetilde L_Z(s)\,(P_Z(J+s)P_Z)^{-1}$, where $\widetilde L_Z(s)=P_Z\widetilde L(s)P_Z$ is the
$14016\times14016$ truncation. In particular the zeros of $\widehat H_{ZZ}$ in $\Omega$, with multiplicities, are
the points $s$ with $-s$ an eigenvalue of the matrix $\widetilde L_Z(0)$, counted with algebraic multiplicity. The same
holds for $R_Z$, with $\mathsf A$ in place of $W\widetilde L_Z(0)W^{-1}$.
\end{lemma}

\begin{proof}
$J+s$ preserves the Fourier index and is upper triangular in $n$, so the range of $P_Z$ is invariant under
$J+s$ and under $Q_0(s)$: $Q_0(s)P_Z=P_ZQ_0(s)P_Z=(P_Z(J+s)P_Z)^{-1}$. Hence
$\widehat H_{ZZ}(s)=P_Z\widetilde L(s)Q_0(s)P_Z=P_Z\widetilde L(s)P_Z\,(P_Z(J+s)P_Z)^{-1}$, and the second factor is
invertible and analytic in $\Re s>-\mu$.
\end{proof}

\emph{The Perron criterion.} To verify that $R+t(H-R)$ is invertible on $\partial\Omega$ for all $t\in[0,1]$, we
use approximate inverses of the diagonal blocks, $V_Z(s)\approx\widehat H_{ZZ}(s)^{-1}$, $V_J\approx\widehat
H_{JJ}(c)^{-1}$ at a fixed centre $c$, and $I$ for the far part, and bound all block entries.

\begin{lemma}[Perron criterion]\label{lem:perron}
Let $X=\bigoplus_iX_i$, $A=R+E$ with $R=\operatorname{diag}(R_i)$, and let $V_i$ be bounded operators with
$\|I-V_iR_i\|\le\rho_i$ and $\|V_iE_{ij}\|\le n_{ij}$. If the spectral radius $\theta$ of the nonnegative matrix
$\operatorname{diag}(\rho_i)+(n_{ij})$ is $<1$, then $R+tE$ is injective for all $t\in[0,1]$; if $R+tE$ is Fredholm of
index zero, it is invertible.
\end{lemma}

\begin{proof}
With $V=\operatorname{diag}(V_i)$, $V(R+tE)=I-\big((I-VR)-tVE\big)$. In the norm
$\max_iw_i^{-1}\|x_i\|$, with $w$ a positive Perron vector of a slightly larger matrix, the operator in
parentheses has norm $\le\theta+\epsilon<1$, since its block norms are entrywise bounded by the matrix above
(and are increasing in $t$). So $V(R+tE)$ is invertible, and $R+tE$ is injective.
\end{proof}

The levels are $Z$, the windows and the far part; this is the \emph{three-level Perron} bound. The entries are
computed as follows (\cref{sec:computation}).
\begin{itemize}[leftmargin=*]
\item \emph{Level $Z$.} $V_Z(s)$ is applied through a Schur decomposition of $\mathsf A$, used
  only as a solver; residuals are certified a posteriori. A bound $t_p\ge\|(\mathsf A+p)^{-1}\|$ is certified by a Loewner-order inequality $t^2MM^*-I\succeq0$ (with preconditioning near the
  B-band root). The coupling $Z\leftarrow$ tail is factored as a low-rank part $F\Phi$ with a certified
  remainder; with the scaling $D_Z=J_Z(c)Q_{ZZ}(s)$, the entry tail$\,\leftarrow Z$ becomes independent of $s$.
\item \emph{Windows.} Each $V_J$ is the exact inverse of the computed block at the centre, and the blocks are
  bounded by block-Jacobi estimates; the dependence on $s$ is bounded by $\rho_J(s)\le\rho_J(c)+d\,K_J/(1-d\,n_{Q,J})$,
  $d=|s-c|$, with constants computed at the centre (and a refined formula for the central windows).
\item \emph{Far part.} Closed forms for $\|Q_0(s)P_{\mathrm{far}}\|$, uniform on discs.
\item \emph{Background and deflation errors.} The difference between $\omegabar$ and the data, and the errors
  $x_j-x_j^A$, enter as global perturbations of all entries.
\end{itemize}

\emph{Discs.} The bound of \cref{lem:perron} is made uniform on a disc around each point $p$ of the contour:
the variation of the level-$Z$ entries is controlled to first order by structured quantities computed by Schur
solves, and to second order by the resolvent bound, using that $\sigma_{\min}(\widetilde L_Z+s)$ is
$1$-Lipschitz in $s$. For each disc the radius is chosen by bisection on a floating-point estimate of $\theta$, with target $0.999$; the
certified rational bounds are then computed for the chosen radius, and may exceed the target slightly, as long as they
stay below $1$. The contour of $\Omega_A$ is
covered by $96$ discs, with $\max\theta\le0.99842$, and that of $\Omega_B$ by the $96$ discs of $\Omega_A$ on its
lower edge and $78$ further discs, with $\max\theta\le0.99908$; the radii range from $0.0028$ near the B-band
root to $0.4995$ at $\Re s\ge2$. Coverage of the four edges by the union of the discs is checked in rational
arithmetic. This proves that $R+t(H-R)$ is invertible on $\partial\Omega_A$ and on $\partial\Omega_B$, hence
\[
  N(\widetilde L_A,\Omega)=N(R,\Omega)\qquad(\Omega=\Omega_A,\Omega_B),
\]
and that $\widetilde L_A$ is invertible on $\partial\Omega_A\cup\partial\Omega_B$. By the remark after
\eqref{eq:brauer-count}, this is \cref{thm:counts}~(iii).

\subsection{The windows have no zeros}\label{sec:counting-windows}

The Perron criterion on the contour gives $N(\widetilde L,\Omega)=N(R,\Omega)=N(\widehat H_{ZZ},\Omega)+
\sum_JN(\widehat H_{JJ},\Omega)$, but says nothing about zeros of the window blocks inside $\Omega$.

\begin{lemma}\label{lem:windows}
For every window $J$, $\widehat H_{JJ}(s)$ is invertible for all $s\in\overline{\Omega_A}\cup\overline{\Omega_B}$.
\end{lemma}

\begin{proof}
Partition each band into $10$ rectangles $\Omega_k$, one for each tail centre $c_k$. On $\Omega_k$ the function
$F_k(s)=I-V_J(c_k)\widehat H_{JJ}(s)$ is analytic, so $\|F_k(s)\|$ is subharmonic and attains its maximum on
$\partial\Omega_k$. On the boundary, $\|F_k\|$ is bounded by the same sensitivity estimates as in the discs,
evaluated on discs around sample points of the boundary. The maximum is $0.1488$ in the A band and $0.1517$ in the B band.
So $\|F_k\|<1$ on $\Omega_k$, and $\widehat H_{JJ}(s)$ is invertible.
\end{proof}

\subsection{The finite count}\label{sec:counting-finite}

It remains to count the eigenvalues of the stored matrix $\mathsf A$ in $-\Omega$. A floating-point Schur decomposition $\mathsf A\approx UTU^*$ gives the matrix $\mathsf B=UTU^{-1}$,
whose eigenvalues are exactly the diagonal entries of $T$ (with $U$ invertible).

\begin{lemma}[Finite count]\label{lem:schur}
Let $\varepsilon_B\ge\|\mathsf A-\mathsf B\|$ and $t(s)\ge\|(\mathsf A+s)^{-1}\|$ on $\partial\Omega$. If
$\sup_{\partial\Omega}t(s)\,\varepsilon_B<1$, then $\mathsf A$ has exactly $\#\{i:-T_{ii}\in\Omega\}$ eigenvalues
in $-\Omega$, with multiplicity.
\end{lemma}

\begin{proof}
$\mathsf A+s+t'(\mathsf B-\mathsf A)$ is invertible on $\partial\Omega$ for $t'\in[0,1]$, by a Neumann series, so the
number of eigenvalues inside does not change along the homotopy.
\end{proof}

We have $\varepsilon_B\le\|\mathsf AU-UT\|_F/\sqrt{1-\|I-U^*U\|_F}$ with $\|\mathsf AU-UT\|_F\le5.97\cdot10^{-5}$ and
$\|I-U^*U\|_F\le1.31\cdot10^{-6}$, and $t(s)\le t_p/(1-rt_p)$ on each disc, so $\sup t\le559.2$ and
$t\,\varepsilon_B\le0.0334$. The diagonal of $T$ has exactly two entries $-T_{ii}$ in $\Omega_A$, at
$0.401025\ldots$ and $0.733180\ldots$, and exactly one in $\Omega_B$, at $0.0414194\ldots+0.5i$; the nearest
entry outside $\Omega_A$ is at distance $0.168$ from its boundary. Hence $N(\widehat H_{ZZ},\Omega_A)=2$ and
$N(\widehat H_{ZZ},\Omega_B)=1$.

\begin{proof}[Proof of \cref{thm:counts}~(ii)]
By \cref{lem:windows,lem:schur,lem:ZZ}, $N(R,\Omega_A)=2$ and $N(R,\Omega_B)=1$. By the Perron verification on the
contours, $N(\widetilde L_A,\Omega)=N(R,\Omega)$, and by~\eqref{eq:brauer-count},
$N(L_A,\Omega_A)=2+1=3$ and $N(L_A,\Omega_B)=1$.
\end{proof}

The positions $0.401$, $0.733$ and $0.0414+0.5i$ of the finite reference are not used in the count; the roots of
$L$ are localized independently in \cref{sec:roots}. The Rouch\'e argument only guarantees that there are three
roots in $\Omega_A$ and one in $\Omega_B$.

\section{Identification of the roots}\label{sec:roots}

By \cref{thm:counts}, $L_A$ has three roots in $\Omega_A$ and one in $\Omega_B$. This section localizes each of
them, proves that each is algebraically simple, and determines the behaviour of its eigenvector with respect to
the constraints. Together with \cref{prop:physmult} this proves the Main Theorem (\cref{sec:roots-proof}).

\subsection{Bordered Newton--Kantorovich}\label{sec:roots-nk}

Fix $c_0$ with $\Re c_0>-\mu$, let $Q_0=(J+c_0)^{-1}$, and let $\ell$ be a bounded functional. Eigenpairs
$(h,\lambda)$ of $L$ with $\ell(h)=1$ correspond, through $h=Q_0y$, to zeros of the bordered map
\[
  F(y,\lambda)=\big(L(\lambda)Q_0\,y,\ \ell(Q_0y)-1\big).
\]
Since $L(\lambda)=L(0)+\lambda$, $F$ is quadratic, with
$DF(y,\lambda)[\eta,\nu]=\big(L(\lambda)Q_0\eta+\nu\,Q_0y,\ \ell(Q_0\eta)\big)$ and a constant second derivative.

\begin{lemma}[Simplicity from the bordered system]\label{lem:bordered}
If $F(y_*,\lambda_*)=0$, then $DF(y_*,\lambda_*)$ is invertible if and only if $\lambda_*$ is an algebraically
simple root of $L$.
\end{lemma}

\begin{proof}
Let $h_*=Q_0y_*$. If $DF[\eta,\nu]=0$ with $\nu\neq0$, then $k=Q_0\eta/\nu$ satisfies $L(\lambda_*)k=-h_*$,
a Jordan chain of length two. If $\nu=0$, then $Q_0\eta\in\ker L(\lambda_*)$ and $\ell(Q_0\eta)=0$; if the
kernel is spanned by $h_*$, where $\ell(h_*)=1$, this forces $\eta=0$. Hence $DF$ is injective if and only if
$\ker L(\lambda_*)=\C h_*$ and there is no Jordan chain, i.e.\ if and only if $\lambda_*$ is algebraically
simple. $DF$ is a finite-rank perturbation of a Fredholm operator of index zero, so injectivity is
equivalent to invertibility.
\end{proof}

We use the Newton--Kantorovich theorem in the following standard form~\cite{NakaoPlumWatanabe2019}. Let
$x_0=(y_0,\lambda_0)$ be an approximate zero and $\mathcal A$ an approximate inverse of $DF(x_0)$, and suppose, in
some norm,
\[
  \|I-\mathcal A\,DF(x_0)\|\le\theta,\qquad \tfrac12\|\mathcal A\,D^2F\|\le Z_2,\qquad \|\mathcal A F(x_0)\|\le Y.
\]
If $\mathcal A$ is injective (here it is invertible by construction) and $(1-\theta)^2>4Z_2Y$, then $T(x)=x-\mathcal AF(x)$ is a contraction of the ball $B(x_0,r)$,
$r=\big(1-\theta-\sqrt{(1-\theta)^2-4Z_2Y}\big)/(2Z_2)$, so $F$ has a unique zero $x_*$ there, and
$\|I-\mathcal A\,DF(x_*)\|<1$, so $DF(x_*)$ is invertible. Uniqueness holds in every ball $B(x_0,\rho)$ with
$\theta+2Z_2\rho<1$.

The norm is a weighted maximum of block norms over the same three levels as in \cref{sec:counting-reference}:
the bordered finite box $Z'$ ($Z$ plus the bordering row and column), the $26$ windows, and the far part. The
bound $\theta$ is the Perron root of the matrix of block bounds, verified in rational arithmetic, and $\mathcal A$
is built from the exact inverse of the computed bordered matrix on $Z'$ and of the computed window blocks. The
centre $\lambda_0$ is an eigenvalue of the truncation $32\times128$ refined by bordered Newton iteration, and
$h_0$ is obtained by inverse iteration on the weighted matrix.

\begin{proposition}[Localization and simplicity]\label{prop:nk}
For each line of the table, there is a unique zero $x_*=(y_*,\lambda_*)$ of $F$ in $B(x_0,r)$, and $\lambda_*$ is an
algebraically simple root of $L_A$ with
\begin{center}
\begin{tabular}{@{}lllllll@{}}
\toprule
root & $\lambda_0$ & $\theta$ & $Z_2$ & $Y$ & $r$ & $|\lambda_*-\lambda_0|\le$ \\
\midrule
$\mu$ & $\mu_A=0.16830707\ldots$ & $0.880172$ & $183.66$ & $6.60\cdot10^{-6}$ & $6.07\cdot10^{-5}$ & $2.82\cdot10^{-5}$ \\
$\sK$ & $0.401024733$ & $0.839937$ & $42.97$ & $5.45\cdot10^{-5}$ & $3.79\cdot10^{-4}$ & $1.72\cdot10^{-4}$ \\
$\sphys$ & $0.7331796596606924$ & $0.805721$ & $10.33$ & $4.19\cdot10^{-6}$ & $2.16\cdot10^{-5}$ & $1.0163\cdot10^{-5}$ \\
$\sB+\frac i2$ & $0.0414194105346554+0.5i$ & $0.91153$ & $155.02$ & $5.97\cdot10^{-6}$ & $7.81\cdot10^{-5}$ & $3.67\cdot10^{-5}$ \\
\bottomrule
\end{tabular}
\end{center}
All quantities are rounded outward from the exact rational values in the certificate files: upward for $\theta$,
$Z_2$, $Y$, $r$ and the error of $\lambda$. The values of $Z_2$ bound $\|\mathcal A\,D^2F\|$ itself, which is
conservative by a factor two. The three roots in the first three lines are the three roots of $L_A$ in $\Omega_A$, and the last one is the root in
$\Omega_B$.
\end{proposition}

\begin{proof}
The Newton--Kantorovich conditions are verified with the quantities in the table (\cref{sec:computation}), and
\cref{lem:bordered} gives simplicity. The intervals of the first three lines are pairwise disjoint and contained
in $\Omega_A$, and each contains a simple root; since $N(L_A,\Omega_A)=3$, these are all the roots in $\Omega_A$.
Likewise for $\Omega_B$, where $N(L_A,\Omega_B)=1$.
\end{proof}

\subsection{The gauge root}\label{sec:roots-gauge}

\begin{proposition}[The gauge root]\label{prop:mu}
$\mu$ is an algebraically simple root of $L_A$, and $\ker L_A(\mu)=\C\,g_*$, with $g_*=(1+Z)\omegabar$.
\end{proposition}

\begin{proof}
$L(\mu)g_*=0$ exactly (\cref{sec:physical-defs}), and $g_*$ lies in the domain, since $g_*-g_A$ is small in norm,
where $g_A=(1+Z)\omega_A$ is computed from the data, and $J+\lambda_0$ is injective on the maximal domain. Normalize
$g_n=g_*/\ell(g_*)$ and let $x_g=(Q_0^{-1}g_n,\mu)$, an exact zero of $F$. We show that $x_g\in B(x_0,r)$ for the first
line of \cref{prop:nk}; then $x_g=x_*$ by uniqueness, so the simple root found there is exactly $\mu$, with
eigenvector $g_*$. From the eigenvalue equation,
\[
  y_g-y_0=-B(g_n-h_0)-L(\lambda_0)h_0+(\lambda_0-\mu)g_n-\delta\mu\,J_1(g_n-h_0),
\]
where $\delta\mu=\mu-\mu_A$ and $J_1$ is the derivative of $J$ in $\mu$. The terms are bounded by
$\|g_*-g_A\|\le2.51\cdot10^{-8}$, which uses the published bounds of~\cite{RT2019} on the difference between
$\omegabar$ and the data together with a bound on $(1+Z)$ of the correction, and by the computed residual
$\|L(\lambda_0)h_0\|\le3.76\cdot10^{-6}$. The result is $\|x_g-x_0\|\le1.82\cdot10^{-5}$, inside the ball $B(x_0,\rho)$, $\rho=3.26\cdot10^{-4}$, on which
$\theta+2Z_2\rho<1$ and the zero of $F$ is unique.
\end{proof}

\subsection{Roots that violate the constraints}\label{sec:roots-witness}

\begin{proposition}[Constraint witnesses]\label{prop:witness}
Let $h_*=Q_0y_*$ be the eigenvector of \cref{prop:nk} at $\sK$, respectively at $\sB+\tfrac i2$. Then, in
$\mathbb T^\sharp(129/128,9/8)$,
\[
  \|\Pi_F\,C(\sK)h_*\|\ge0.16975,\qquad \|\Pi_F\,C(\sB+\tfrac i2)h_*\|\ge0.35469,
\]
where $\Pi_F$ is the projection onto the box $F=\{|m|<12,\ n<48\}$. In particular both eigenvectors violate the
constraints, and so does the eigenvector of $L$ at $\sB$ (\cref{lem:sectorB}).
\end{proposition}

\begin{proof}
At the centre, $c_{\mathrm{ref}}=\|\Pi_FC(\lambda_0)h_0\|$ is computed rigorously on the box. The difference
$\Pi_F\big(C(\lambda_*)h_*-C(\lambda_0)h_0\big)$ is bounded level by level using the Newton--Kantorovich ball:
the level $Z'$ through $\|\Pi_FCQ_0P_Z\|$, the windows through the decay of the columns (tail columns start at
$n=128$ and reach the box only through the factor $(5/4)^{-80}$), the far part by closed forms, and the variation
of $\lambda$ through $C_1$. For $\sK$: $c_{\mathrm{ref}}\ge0.171739$, and the losses total at most $1.981\cdot10^{-3}$, which gives
$0.171739-0.001981=0.169758$. For $\sB$: $c_{\mathrm{ref}}\ge0.355188$, and the losses total at most
$4.954\cdot10^{-4}$, which gives $0.355188-0.0004954=0.3546926$.
\end{proof}

\subsection{The physical root: injectivity of $K$}\label{sec:roots-K}

The eigenvector at $\sphys$ cannot be shown to satisfy the constraints by a lower bound; instead we show that it
cannot violate them. By~\eqref{eq:KCML}, if $L(s)h=0$ then $K(s)C(s)h=0$, so $C(s)h\neq0$ is only possible at a
root of $K$.

\begin{proposition}[Injectivity of $K$]\label{prop:Kinj}
In $\mathbb T^\sharp(129/128,9/8)$, $K(s)$ is injective for every $s$ with $\Re s\ge0$ and $0\le\Im s\le\tfrac14$,
except at one point $s=\sK'$, with $|\sK'-0.401|<0.05$, which is an algebraically simple root of $K$.
\end{proposition}

\begin{proof}
The region is covered by three certificates.
\begin{itemize}[leftmargin=*]
\item $\Re s>1.765$: \cref{thm:counts}~(i).
\item $\{0\le\Re s\le1.765,\ 0\le\Im s\le\tfrac14\}$ minus the open disc $|s-0.401|<0.05$: a cover by $154$
  tiles. On each tile, injectivity of $K(s)$ follows from \cref{lem:perron} applied to $K(s)Q_0(c)$, with four
  levels (a finite box, two shells and the far part) and a Perron root $\le0.99612$; the cover is checked in rational
  arithmetic by a quadtree.
\item The disc $|s-0.401|\le0.05$: an operator Rouch\'e argument on its boundary circle, as in
  \cref{sec:counting-reference}, with reference $\operatorname{diag}(\widehat H_{ZZ}(s),I,I,I)$ on a box
  $Z=20\times80$ and Perron roots $\le0.8126$ on $16$ tiles covering the circle (the signed norm is not symmetric
  under conjugation, so the whole circle is covered). The finite reference is counted by a bordered scalar
  Rouch\'e argument: the truncation has exactly one eigenvalue in the disc, simple, near $0.4010247311$.
\end{itemize}
\end{proof}

\begin{corollary}\label{cor:sK}
$\sK'=\sK$, and the eigenvector of $L$ at $\sphys$ satisfies $C(\sphys)h=0$.
\end{corollary}

\begin{proof}
The eigenvectors of \cref{prop:nk} are computed in the torus space; since $|\sK|,|\sphys|\le3.1$, they lie in the
domain of $L$ in $\Yp$ by \cref{lem:Lreal}, so \cref{lem:KCML} applies to them with $(a,b)=(65/64,5/4)$ and
$(a',b')=(129/128,9/8)$. By \cref{prop:witness}, $C(\sK)h_*\neq0$, and $K(\sK)C(\sK)h_*=M(\sK)L(\sK)h_*=0$, so $\sK$ is
a root of $K$ in the region of \cref{prop:Kinj}; hence $\sK=\sK'$. Since $\sphys$ is real and $|\sphys-0.401|>0.33$,
$K(\sphys)$ is injective, so $C(\sphys)h=0$ for the eigenvector $h$. Here $C(\sphys)h$ lies in the domain of $K$ in $Y^\sharp(129/128,9/8)\subset\mathbb T^\sharp(129/128,9/8)$, where
\cref{prop:Kinj} applies (\cref{app:KCML}).
\end{proof}

\subsection{Reality}\label{sec:roots-reality}

\begin{proposition}\label{prop:real}
$\sK$, $\sphys$ are real, and $\Im(\sB+\tfrac i2)=\tfrac12$, i.e.\ the root $\sB$ of $L$ is real.
\end{proposition}

\begin{proof}
In the realization $\Yp$, the conjugation $h_{mn}\mapsto\overline{h_{-m,-n}}$ is an isometry, and
$\overline{L_A(\bar s)}=L_A(s)$ in the sense of \cref{sec:setting-sectors}; so the roots of $L_A$ in $\Yp$, with
multiplicities, are symmetric under $s\mapsto\bar s$, and those in the B band under $s\mapsto\bar s+i$. By
\cref{lem:Lreal} the same holds in the torus realization for $|s|\le3.1$. In $\Omega_A$ the roots are $\mu,\sK,\sphys$,
with $\sK$ and $\sphys$ in discs of radius $<2\cdot10^{-4}$ about real centres; the conjugate of $\sK$ is a root in the
same disc, hence equal to $\sK$, and likewise for $\sphys$. In $\Omega_B$ there is a single root, so it is fixed by
$s\mapsto\bar s+i$, whose fixed set is $\Im s=\tfrac12$.
\end{proof}

\subsection{Proof of the Main Theorem}\label{sec:roots-proof}

\emph{Part~(1).} By \cref{cor:four}, modulo $\tfrac i2\Z$ there are four roots in $\Re s>0$, with multiplicity. By
\cref{prop:nk,prop:mu,prop:real}, they are $\mu$, $\sK$, $\sphys$ (sector A) and $\sB$ (sector B, appearing as
$\sB+\tfrac i2$ for $L_A$), all real and algebraically simple, with the stated enclosures; the interval for
$\sphys$ is $\lambda_0\pm1.0163\cdot10^{-5}$, rounded outward to eight digits. The eigenvectors at $\sK$ and $\sB$ violate the constraints
(\cref{prop:witness}); the eigenvector at $\mu$ is $g_*$, i.e.\ $e^{\mu\tau}g_*$ is the gauge mode $G_*$
(\cref{prop:mu}); the eigenvector at $\sphys$ satisfies the constraints (\cref{cor:sK}). The statement on the
imaginary axis is \cref{cor:four}, with eigenvector $\partial_\tau\omegabar$ at $s=0$.

\emph{Part~(2).} All four roots are simple, so \cref{prop:physmult} applies in its simple form:
\begin{itemize}[leftmargin=*]
\item at $\sK$ and $\sB$, $C(s)h\neq0$, so the physical multiplicity is $0$;
\item at $\mu$, $h=g_*$ spans $\mathcal G_\mu$, so it is $0$, and $1$ in the cone-fixed setting;
\item at $\sphys$, $C(\sphys)h=0$ and $\sphys\not\equiv\mu\pmod{\tfrac i2\Z}$, so $\mathcal G_{\sphys}=0$ and it is $1$.
\end{itemize}
By \cref{thm:correspondence}, the physical Floquet modes with $\Re s>0$, modulo gauge, form a one-dimensional space,
with exponent $\sphys$; since $\sphys$ is simple, there are no generalized modes. In the cone-fixed setting the
count is two, the additional mode being $\Lie_{X_0}(\gbar,\phibar)$, which by \cref{lem:cone} is tangent to the
translations of the singular point. \qed

\section{Rigorous computation and reproducibility}\label{sec:computation}

This section describes what the computer-assisted parts of the proof consist of, what has to be trusted, and how
a reader can check them.

\subsection{Arithmetic and trust base}\label{sec:computation-arith}

Three kinds of computation enter the certificates.
\begin{enumerate}[label=(\alph*),leftmargin=*]
\item \emph{Exact rational arithmetic} (Python's \texttt{fractions}): the final inequalities of every
  certificate, namely the Perron conditions $Nv\le\theta v$ with $\theta<1$ of
  \cref{lem:perron}, the Newton--Kantorovich conditions of \cref{sec:roots-nk}, the covers of the contours and of
  the region of \cref{prop:Kinj}, the deflation polynomial $q$, and the fact F1.
\item \emph{Ball arithmetic} with Arb~\cite{Arb}, through \texttt{python-flint}: the symbol bounds of
  \cref{lem:F3}, and the slack between grid points.
\item \emph{A posteriori bounds on floating-point computations.} Large linear algebra (dimension $14016$ for the
  finite reference, windows of dimension up to $11200$) is done in IEEE double precision. Floating-point results
  are used only as approximate inverses, solvers or test vectors, and every quantity that enters a certificate is
  bounded rigorously: residuals are computed with entrywise rounding bounds of the standard form
  $\gamma_n\||M||X|\|$~\cite{Higham2002}; norms of inverses are bounded by verified Loewner inequalities
  $t^2MM^*-I\succeq0$, whose positive definiteness is verified by Rump's method~\cite{Rump2006} applied to the
  realification; and the coupling factors are certified as in \cref{sec:counting-reference}.
\end{enumerate}
The trust base is therefore: the correctness of the bounds published in~\cite{RT2019}, which we use as stated;
IEEE~754 arithmetic with rounding to nearest, as assumed by the a posteriori bounds; Arb; and the Python code of the
verification program, which is short for each check and is available for inspection. No part of the final chain
is formalized in a proof assistant.

Because the floating-point entries depend on the BLAS library and on the number of threads, re-running a
certificate on another machine reproduces the rational $\theta$ only up to $\sim10^{-16}$, not exactly; the
criterion for a successful reproduction is $\theta<1$ with a difference at the scale of rounding.

\subsection{The certificates}\label{sec:computation-certs}

\Cref{tab:certs} lists the computed certificates, the results they establish, and their approximate cost. The
computations were run on five commodity computers, with four to six cores and $7$--$31$\,GB of memory; the whole
chain takes a few days on three or four such machines. The largest single computations are the Schur
decomposition of the finite reference and the Loewner bounds on its resolvent (hours per contour centre), and the
tails with $26$ dense windows (three to six hours per centre).

\begin{table}[ht]
\centering
\small
\begin{tabular}{@{}p{0.29\textwidth}p{0.35\textwidth}p{0.27\textwidth}@{}}
\toprule
certificate & establishes & cost \\
\midrule
symbols of $B$ and $B^\sharp$ (Arb, grids $2048\times4096$ and $1024\times2048$) & \cref{lem:F3}; \cref{thm:counts}~(i) & hours, 4 processes \\
radius recovery and transfer & \cref{lem:Lreal,lem:recovery} & minutes \\
finite reference: matrix, Schur, residual & \cref{lem:schur} & hours \\
level $Z$ on $174$ discs (Loewner, Schur solves) & \cref{thm:counts}~(ii),(iii) & $\sim$9 min per point \\
low-rank coupling factors ($13$ centres) & \cref{sec:counting-reference} & $\sim$1 h per centre \\
tails, $26$ windows ($13$ centres) & \cref{sec:counting-reference} & 3--6 h per centre \\
windows inside the bands & \cref{lem:windows} & hours \\
bordered NK at $\mu,\sK,\sphys,\sB$ & \cref{prop:nk,prop:mu} & 3--6 h each \\
constraint witnesses & \cref{prop:witness} & minutes \\
$154$ tiles and Rouch\'e for $K$ & \cref{prop:Kinj} & $\sim$1 day, 3 machines \\
F1 & \cref{lem:F1} & seconds \\
\bottomrule
\end{tabular}
\caption{Computed certificates. File names and sha256 hashes are listed in \cref{app:certificates}.}
\label{tab:certs}
\end{table}

\subsection{The verification program}\label{sec:computation-verify}

The repository contains a verification program with two levels.

\emph{Level 1} re-checks, in exact rational arithmetic and with the code of the repository, the final inequalities
of the certificates from the stored certificate files (small JSON files):
\begin{itemize}[leftmargin=*]
\item the three-level Perron bound at each of the $175$ points of the $22$ disc files of the two bands ($174$ discs of
  positive radius, and one entry of radius zero that is not used in the cover);
\item the covers of $\partial\Omega_A$ and $\partial\Omega_B$;
\item the finite counts;
\item the window bounds;
\item the four Newton--Kantorovich certificates, the witnesses and the identification of the gauge root;
\item the $154$ tiles and their cover of the region of \cref{prop:Kinj} outside the disc $|s-0.401|<0.05$ of the
  Rouch\'e argument for $K$;
\item the deflation polynomial $q$ of \cref{sec:counting-rouche};
\item the radius recovery and transfer bounds, the bound $\beta_+$ on the background part, and F1.
\end{itemize}
Some certificates enter level 1 only through their stored values, and are recomputed only at level 2: the symbol
bounds and the free parts (level 1 recombines the stored bands), and the Perron bounds of the $16$ tiles of the
Rouch\'e argument for $K$.
It takes about two hours on four cores. All checks pass: the worst Perron root on the discs is $0.999074$, with
differences to the stored values $\le7.8\cdot10^{-16}$, and the Newton--Kantorovich data, witnesses and tiles are
reproduced identically.

\emph{Level 2} recomputes the intermediate certificates from the data of~\cite{RT2019}. As a sample, we recomputed
from scratch on a different machine:
\begin{itemize}[leftmargin=*]
\item the reference matrix, identical bit for bit to the stored one;
\item the Schur residual bound, identical;
\item the tail of the contour centre $3.0$ ($26$ windows and the far level, $5.3$ hours), identical in every field.
\end{itemize}

\subsection{Data and code availability}\label{sec:computation-data}

The code, the certificate files, the verification program and a reproduction guide are available at
\url{https://github.com/edivaldo-amaral/choptuik-espectro} (MIT licence)~\cite{Repo2026}. The large data files, namely the reference
matrix, its Schur factors and the coupling factors ($16$\,GiB), are deposited at Zenodo, DOI
\href{https://doi.org/10.5281/zenodo.23222363}{10.5281/zenodo.23222363} (CC BY 4.0)~\cite{Zenodo2026}, with sha256 hashes that are
also recorded in the certificate files. The data of~\cite{RT2019} are not redistributed; a script downloads them
from the arXiv source of~\cite{RT2019} and checks their sha256 hashes.
\todo{Repository and Zenodo record become public with the preprint.}

\subsection{Independent AI-assisted checks}\label{sec:computation-reviews}

Before being used, the certificates and the lemmas were checked by AI agents independent of those that produced the
work, with the exceptions listed below. These checks are not reviews by human experts. The protocol was designed so
that a check produces verifiable evidence rather than an opinion:
\begin{itemize}[leftmargin=*]
\item the checking agent ran on a different model from the one that produced the work, with a clean context;
\item it was given the precisely stated claims, the derivations under review, and the code and data, but not
  the author's conclusions or earlier audits;
\item it had to deliver one artifact per claim: an independent derivation, an independent implementation, or a
  direct computation of the true quantity to compare with the bound;
\item a snapshot with sha256 manifest was taken before each check and compared after it.
\end{itemize}
There were ten such checks:
\begin{itemize}[leftmargin=*]
\item the tiles of \cref{prop:Kinj} (an earlier check of them with another tool was superficial and is not
  counted);
\item the certificates at $\sK$ (with the Rouch\'e argument for $K$) and at $\mu$;
\item the Rouch\'e arguments of the two bands, including the certificate at $\sB$;
\item the structural lemmas and the chain of the proof;
\item the gauge lemma, the forward reconstruction and the converse;
\item the modelling assumptions.
\end{itemize}
Two parts have not been checked on their own. The Newton--Kantorovich certificate at $\sphys$ uses the same code as
the one at $\sB$, which was checked. The symbol computations of \cref{lem:F3} were examined only in their use, not
recomputed. None of the checks overturned a result. All of them found gaps, mostly of traceability or of statement,
and some of derivation: for instance, one found that an early version of the tiles used the dyadic approximation of
$\mu$ instead of the true value, and another that the deflation error outside the finite box had been omitted. All
were corrected before the results were used; the reports and artifacts are in the repository.

\subsection{Use of AI assistants}\label{sec:computation-ai}

This work was done with extensive use of AI assistants, between September and October 2026. We describe their role
as required by the journal's policy.

\emph{Scope.} The assistants wrote most of the code, carried out most of the derivations and computations, prepared
the documentation, and drafted this manuscript. Their contribution is substantial and not limited to language
editing.

\emph{The author's role.} The author posed the problem, the unstable spectrum of the solution of~\cite{RT2019}, which
its authors left open; directed the project and chose among the approaches; ran the computations on the author's own
and on borrowed machines; checked the arguments and the results; and takes full responsibility for the content of
this paper.

\emph{Tools.} The main assistant was Claude (Anthropic), used through the Claude Code environment, with the models
\texttt{claude-opus-5} and later \texttt{claude-opus-5-5}. OpenAI's Codex agent was used in the first exploratory
session and during a three-day interval, with the models \texttt{gpt-5.6-luna}, \texttt{gpt-5.6-sol},
\texttt{gpt-6-astra} and \texttt{gpt-reserve}; everything produced in that interval was audited and re-derived or
checked before use. The checks of \cref{sec:computation-reviews} were done by Claude agents running
\texttt{claude-fable-5-1}, a model different from the one that produced the work. At the author's request, a draft of the
manuscript was also read by ChatGPT (OpenAI, model Sol 6.1); its comments were checked against the certificates and
the sources before being incorporated. The model identifiers are those recorded in the session logs; a dated account of the sessions is in the file \texttt{paper/AI\_USE.md} of the
repository.

\emph{Limits of the verification.} The analytic arguments are given in this paper; some routine estimates appear in
sketched form, with their constants checked by code in the repository (\cref{app:lemmas}). Every computed bound is
produced by code that can be inspected and re-run from the certificate files and data, and all of them have been
re-checked by the verification program (level~1 of \cref{sec:computation-verify}); a sample of the intermediate
certificates has been recomputed from the data of~\cite{RT2019} (level~2). The lemmas and most certificates were also
examined by the checks of \cref{sec:computation-reviews}, with the exceptions listed there. None of this replaces
peer review; the aim is that trust in the result should not depend on the process that produced it.

\section{Discussion}\label{sec:discussion}

\subsection{What the theorem says about critical collapse}

In the standard picture of critical collapse~\cite{GundlachMartinGarcia2007}, initial data near the threshold of
black hole formation evolve for a long time close to the critical solution, and the deviation is dominated by its
single growing mode, $e^{\lambda T}$. If the initial data depend on a parameter $p$, the black hole mass then
scales like $(p-p_*)^\gamma$ with $\gamma=1/\lambda$, modulated by a periodic wiggle in the discretely self-similar
case. The Main Theorem supplies the spectral input of this picture for spherically symmetric scalar field
collapse: there is exactly one unstable mode, it is simple, and $\lambda\in[2.674040,\,2.674115]$. Turning the
picture into a theorem needs further ingredients, among them the following; the regularity question of
\cref{sec:discussion-regularity} and the neutral modes of \cref{rem:neutral} would also have to be addressed.
\begin{enumerate}[label=(\alph*),leftmargin=*]
\item \emph{Linear stability in the evolution sense.} Absence of other roots of $L$ in $\Re s>0$ does not by
  itself give decay of the linearized flow on a complement of the unstable and gauge modes: the operator is not
  self-adjoint, and one needs resolvent bounds along vertical lines, or a completeness statement for the Floquet
  modes, rather than the location of eigenvalues alone. Reiterer has given an abstract framework for finite
  codimension stability of time-periodic symmetric hyperbolic equations with a compact resolvent sufficiently far to the right in the complex
  plane, motivated by discretely self-similar spacetimes~\cite{Reiterer2015}; whether its hypotheses hold for the
  present formulation is not addressed here.
\item \emph{Nonlinear stability:} a codimension-one stable manifold of the critical solution.
\item \emph{Global aspects.} The solution of~\cite{RT2019} is constructed on a neighbourhood of the past light cone
  of the singularity; the threshold picture concerns asymptotically flat data.
\end{enumerate}
The numerical agreement of $1/\lambda$ with the measured exponent $\gamma\approx0.374$ is therefore a consistency
check, not a consequence.

\subsection{Nonspherical perturbations}

Reiterer and Trubowitz describe the numerical evidence on general perturbations as inconclusive~\cite{RT2019},
citing the linear analysis of Mart\'in-Garc\'ia and Gundlach~\cite{MartinGarciaGundlach1999}, who found all
nonspherical perturbations to decay, and the axisymmetric nonlinear evolutions
of~\cite{ChoptuikHirschmannLieblingPretorius2003}. The delicate sector is the polar (even) sector with $\ell=2$,
whose least damped mode was found at $\kappa\Delta\approx-0.07$, close to the imaginary axis.

As a first step toward extending the Main Theorem, we have written the Gerlach--Sengupta
equations~\cite{GerlachSengupta1979} on the background of~\cite{RT2019}, in the variables $\omega$ and $\mu$, and
solved them numerically by spectral collocation, keeping the light cone inside the domain and filtering spurious
solutions by regularity and by the constraints. This computation is \emph{not} rigorous and is not part of the
theorem. In both sectors, for $\ell=2$ and $\ell=3$, it is consistent with the values
of~\cite{MartinGarciaGundlach1999} within the precision of their finite-difference computation; for instance
$\kappa\Delta\approx-2.320$ against their $-2.30$ for the axial mode with $\ell=2$. For the polar mode with $\ell=2$
it gives
\[
  s\approx-0.004764\pm0.147345\,i,\qquad \kappa\Delta=4\pi\Re s\approx-0.0599,
\]
against their $-0.07$, and stable under refinement of our discretization to about $2\cdot10^{-5}$, which is
numerical evidence and not a certified bound: the mode is damped, but at a distance of only $0.005$ from the
axis. A rigorous count for $\ell\ge1$ appears feasible with the methods of this paper, using a deflation and a
Newton--Kantorovich enclosure for this mode, the constraint structure in the certificate, and a bound uniform in
large $\ell$. Together with the Main Theorem it would give exactly one unstable mode for general linear perturbations,
in the spectral sense.

\subsection{Regularity}\label{sec:discussion-regularity}

The class of physical modes in \cref{def:physical-mode} requires analyticity in $\xi$ across the light cone and
at the axis, with a complex neighbourhood that may be arbitrarily thin. Whether every Floquet mode with
$\Re s>0$ that is merely smooth in $\xi$ is analytic is an open regularity question for a hyperbolic system that is
periodic in $\tau$. Our radius recovery (\cref{lem:recovery}) works for every weight $\kappa_2'>1$, but its bound on
the free part degenerates like $1/(\kappa_2'-1)$, so it does not reach $\kappa_2'=1$.

\subsection{On the method}

What made the count feasible on ordinary computers was a combination of choices of formulation.
\begin{itemize}[leftmargin=*]
\item \emph{A finite reference instead of a majorant of the exterior.} An earlier approach bounded the resolvent of
  the finite section and the operator outside it separately; it required truncations and constants out of reach
  without high-performance computing. Comparing $L$ with a block-diagonal reference whose zeros are those of a
  finite matrix needs only bounds along the contour.
\item \emph{A space in which the background is pointwise.} On the torus space the numerical range of the background
  reduces to the symbol, which gave the exclusion bound $2.93$, against $414$ from block-by-block majorants.
\item \emph{The Perron criterion} for block matrices of norms, instead of the spectral norm of such a matrix.
\item \emph{Deflation of the gauge roots,} which lie on or near the contour.
\item \emph{Constraints through $K$.} In a floating-point diagnostic on the truncation $12\times36$, the maximum
  over the boundary of $[1/8,1]\times[-1/4,1/4]$ of the weighted norm of $J(\,\cdot\,+s)^{-1}$ was about $1800$ for
  $L$ and $14$ for $K$, each in the space where its certificate is formulated. The physical classification is therefore obtained by
  proving $K$ injective, rather than by bounding the constraints on root spaces of $L$.
\end{itemize}
Several other routes were tried and closed, with measured lower bounds on their cost. They are documented in the
repository, for readers who might attempt similar computations.


\section*{Acknowledgements}

The author thanks Jo\~ao Rondon and Andr\'e Leonam, the author's sons, for lending the idle hours of their
computers, and the Departamento de Engenharia El\'etrica e Automa\c{c}\~ao (DEEA) of the Instituto Federal de Mato
Grosso for the use of two computers.

\section*{Statements and declarations}

\emph{Funding.} This work received no specific funding.

\emph{Competing interests.} The author has no competing interests to declare.

\emph{Data and code availability.} The code, the certificate files and the verification program are available
in the repository~\cite{Repo2026}; the large data files are deposited at Zenodo~\cite{Zenodo2026}
(\cref{sec:computation-data}).

\emph{Use of AI tools.} AI assistants were used extensively in this work; their role, the models used and the
verification procedures are described in \cref{sec:computation-ai,sec:computation-reviews}.


\appendix
\section{Proofs of the structural lemmas}\label{app:lemmas}

\subsection{Sectors, shift and conjugation (\cref{lem:sectorB})}\label{app:sectors}

Let $(Uh)_{mn}=h_{m-1,n}$, i.e.\ $Uh=e^{i\tau/2}h$. In $\mathbb T(a,b)$ the weights $a^{2m}$ give
$\|Uh\|=a\|h\|$, and in $Y(a,b)$ we have $a^{-1}\|h\|\le\|Uh\|\le a\|h\|$; so $U$ is an isomorphism of each space,
an isometry up to a constant factor. It commutes with convolution by a field (a sum over $m$), with the reflection
$P$, with multiplication by $\xi$ and with $R_\xi$, all of which act in $n$ only or are translation invariant in
$m$; and $U^{-1}\partial_\tau U=\partial_\tau+\tfrac i2$. The free part $J$ therefore satisfies
$U^{-1}(J+s)U=J+s+\tfrac i2$, $U$ maps the domain of $J$ onto itself, and $U^{-1}L(s)U=L(s+\tfrac i2)$,
$U^{-1}C(s)U=C(s+\tfrac i2)$. Since $U$ changes the parity of $m$ in every component, it exchanges $X_A$ and $X_B$.
Conjugation by $U$ maps kernels and Jordan chains to kernels and Jordan chains, and constraint-satisfying ones to
constraint-satisfying ones. This proves \cref{lem:sectorB}.

\emph{Conjugation.} Let $(\mathcal Ch)_{mn}=\overline{h_{-m,n}}$, which corresponds to complex conjugation of the
function $h(\tau,\xi)$ for real $\tau,\xi$ (recall $h_{m,-n}=h_{mn}$). The background is real, so
$\mathcal C L(s)\mathcal C=L(\bar s)$, and $\mathcal C$ preserves both sectors. In $Y(a,b)$, $\mathcal C$ is an
isometry, so the roots of $L_A$ and of $L_B$ in $Y(a,b)$, with their multiplicities, are invariant under
$s\mapsto\bar s$. (In the torus space $\mathcal C$ is not an isometry, since the weight $a^{2m}$ is not even in $m$;
this is why the reality argument of \cref{prop:real} goes through $\Yp$ and \cref{lem:Lreal}.) For the B band,
$\sigma(L_A)\cap\{\tfrac14<\Im s<\tfrac34\}=\sigma(L_B)+\tfrac i2$, and the invariance of $\sigma(L_B)$ under
conjugation becomes invariance under $s\mapsto\overline{s-\tfrac i2}+\tfrac i2=\bar s+i$.

\subsection{The Fredholm property (\cref{prop:fredholm})}\label{app:fredholm}

\emph{(i) The free part.} $J+z$ acts on each Fourier mode $m$ and each component separately, as an upper triangular
operator in $n$ with diagonal $\mu(n+1)+z+im/2$ and first-order off-diagonal structure; the explicit inverses
of~\cite[\S3]{RT2019} are bounded in $Y(a,b)$, with a bound of the form $C/(\mu+z)$, and in the torus space they are
bounded mode by mode by closed forms (in the repository). The norm of $Q(I-G)$, where $G$ is the projection
onto a box $\{|m|<M,\,n<N\}$, tends to zero as $M,N\to\infty$; for instance, with $r=1/b$,
\[
  \|Q(I-G)\|\le\max\Big(\frac{3(1+2\mu_{\max})}{M(1-r)},\ \frac{3(1+4/(1-r))}{2\mu_{\min}(N+1)}\Big)
\]
in $Y(1,b)$, where $\mu_{\min}\le\mu\le\mu_{\max}$. Hence $Q$ is a norm limit of finite-rank operators and is compact.
The maximal domain $\Dom_{\max}=\{u:Ju\in Y\}$, with $Ju$ computed termwise, coincides with the closure
$\Dom_{\min}$ of the trigonometric polynomials in the graph norm: for $z>0$, $Q(z)$ maps $Y$ into $\Dom_{\min}$ and
inverts $J+z$ there; for $u\in\Dom_{\max}$, $w=u-Q(z)(J+z)u$ satisfies $(J+z)w=0$ as an analytic function, and in
each Fourier mode the homogeneous solutions are multiples of $(1+\xi)^{-1-(z+im/2)/\mu}$ (components with
$(1+\xi)\partial_\xi$) or $\xi^{-1-(z+im/2)/\mu}$ (the component with $\xi\partial_\xi$), which are singular at
$\xi=-1$ or $\xi=0$, inside the Bernstein ellipse of parameter $b>1$. So $w=0$.

\emph{(ii) The background.} $B=\mu S\Gamma_1+2S\Xi\Gamma_2(\omegabar,\cdot)$ is a sum of reflections, the regularized
division $R_\xi$, multiplication by $\xi$ and convolution with the coefficients of $\omegabar$. The latter lie in
$\ell^1$ with weights $(65/64)^{|m|}(5/4)^{|n|}$, so convolution is bounded in every $Y(a,b)$ and $\mathbb T(a,b)$
with $a\le65/64$, $b\le5/4$; $R_\xi$ is bounded for $b>1$~\cite[\S4]{RT2019}.

\emph{(iii)} $L(s)Q=I+(B+s-z_0)Q$ is the identity plus a compact operator, analytic (affine) in $s$, hence Fredholm of
index zero. For real $s$ large, $L(s)=(J+s)\big(I+(J+s)^{-1}B\big)$ is invertible in $Y(a,b)$ by a Neumann series, since
$\|(J+s)^{-1}\|\le C/(\mu+s)$; in the torus space, $L(s)$ is invertible for $s>2.9261$ by \cref{thm:counts}~(i). By the analytic Fredholm theorem the roots are isolated and of finite algebraic multiplicity. The same
argument applies to $K(s)=J^\sharp+B^\sharp+s$ in the sharp spaces.

\subsection{Radius recovery (\cref{lem:Lreal,lem:recovery})}\label{app:recovery}

Let $Y_s\subset Y_w$ be a strong and a weak realization (for \cref{lem:Lreal}, $\Yp\subset\mathbb T(65/64,5/4)$;
for \cref{lem:recovery}, $\Yp\subset Y(\kappa_1',\kappa_2')$), with the same operator $Q=J^{-1}$, which preserves
boxes. Write $H(s)=L(s)Q=I+(B+s)Q$, let $G$ be the projection onto a finite box and $T=I-G$, and suppose
\[
  \|T(H(s)-I)T\|_{Y_s}<1,\qquad \|T(H(s)-I)T\|_{Y_w}<1 .
\]
\emph{Kernels.} Let $H(s)x=0$ with $x\in Y_w$. Then $Gx$ is a finite vector, hence in $Y_s$, and
$THT\cdot Tx=-THGx$, whose right side is in $Y_s$ because $H$ is bounded on $Y_s$. In $Y_s$ the operator $THT$ on the
range of $T$ is invertible by a Neumann series, so there is $y\in Y_s$ with $THT\,y=-THGx$. The same equation holds
in $Y_w$, where its solution is unique; so $Tx=y\in Y_s$, and $h=Qx$ lies in the domain in $Y_s$. The converse
inclusion is the inclusion of spaces. \emph{Chains.} If $L(s)h_j=-h_{j-1}$ with $h_{j-1}$ already known to be in the
strong domain, the same argument applies to $x_j=Jh_j$ with the additional strong right-hand side $-Th_{j-1}$.

\emph{Constants for \cref{lem:Lreal}.} Box $G=\{|m|<1024,\,n<4096\}$, $|s|\le3.1$. In the torus space,
$\|T(H-I)T\|\le(\|B\|+|s|)\|QT\|\le(3.9642+3.1)\cdot0.0525<0.371$, with $\|B\|$ from the symbol (\cref{lem:F3}) and
$\|QT\|\le0.0525$ from closed forms: the tail $n\ge4096$ in the modes $|m|<400$, and the bound
$(1+2c_w/(\kappa_2-1))/(|m|/2)$ for the whole mode when $|m|\ge400$. In $\Yp$, with component weights
$\eta=(916,4096,243,1233)$ and the bound $\beta_+\le5191/500$ on the background part, computed in exact
arithmetic from the data and the published ball of~\cite{RT2019} (\cref{app:certificates}), the same quantity is
$\le0.2666$. The weights satisfy $\eta_i\ge243$, so $\|x\|_{\mathbb T}\le\|x\|_{\Yp}/243$.

\emph{Constants for \cref{lem:recovery}.} In the weak space $Y(\kappa_1',\kappa_2')$ the background part is bounded by
$\beta(\kappa_2')=\beta_+\,D(\kappa_2')/D(5/4)$ with $D(k)=2k/(k^2-1)$, the factor coming from the bound on $R_\xi$; the
Fourier weight $\kappa_1'\in[1,65/64]$ does not increase the convolution norms, and $\|QT\|$ does not depend on it
because $Q$ preserves $m$. Since $\|QT\|\to0$ for every $\kappa_2'>1$, for every $S$ and $\kappa_2'$ there is a box with
$(\beta(\kappa_2')+S)\|QT\|<1$; explicit boxes are computed, e.g.\ $2^{17}\times2^{20}$ for $S=200$ and
$\kappa_2'=1+2^{-6}$. The rational verification is in the repository (\cref{app:certificates}).

\subsection{Propagation of the constraints (\cref{lem:KCML})}\label{app:KCML}

Let $O=\tfrac12(1-P)$, $E=\tfrac12(1+P)$, $X$ multiplication by $\xi$, and for a residual $q=\Omega^+(\mu,w)$ with
components $(q_1,q_2,q_{2\sharp},q_3,q_4)$ set $F=Sq$, $c=S^\sharp q=(Oq_1,-Pq_{2\sharp})$ and $e=-Pq$. For residuals of
the system, the selected components are divisible by $\xi$ (the principal and quadratic terms carry a factor
$\xi$, and the selected components of $\Gamma_1$ are odd), so that
\[
  e=Ic+RF,\qquad Ic=(c_1,0,c_2,0,0),\qquad Rf=(XPf_1,XPf_2,0,XPf_3,XPf_4),
\]
using $PX=-XP$ and $Pc_1=-c_1$. Reiterer and Trubowitz prove the differential identity of \cref{eq:sharp} in an
\emph{off-shell} form~\cite[\S2]{RT2019}: with $J_s=\partial_\tau+s+\mu((1+\xi)\partial_\xi+1)$ and $\mathcal G(w,e)$
the bilinear map of their identity,
\[
  T_s(w)e:=X\big(OJ_se_1,\ J_se_{2\sharp}-PJ_se_2\big)+\mu\big(Ee_1,\ e_1+e_2+Pe_2-2Pe_3\big)+X\mathcal G(w,e)
\]
satisfies $T_0(w)e(w)=0$ for every admissible $w$, whether or not $e(w)=0$. The only dependence of $T$ on $w$ is
through $\mathcal G$, so linearizing at $w$ in the direction $e^{s\tau}h$ gives $T_s(w)e'+X\mathcal G(h,e(w))=0$,
where $e'$ is the linearized residual. With $F'=L(s)h$, $c'=C(s)h$, $e'=Ic'+RF'$, and
\[
  K(s)c:=R_\xi T_s(w)Ic,\qquad M(s)f:=-R_\xi T_s(w)Rf,
\]
and $R_\xi X=I$, we obtain $K(s)C(s)h=M(s)L(s)h-R_\xi X\mathcal G(h,e(w))$. At the exact background $e(\omegabar)=0$,
which is~\eqref{eq:KCML}. The explicit forms of $K$ and $M$ are in \cref{app:formulas}; $K$ has the principal part
$\operatorname{diag}(\partial_\tau+s+\mu(\xi\partial_\xi+1),\,J_s)$ of~\eqref{eq:sharp}.

\emph{Domains.} $M(s)$ contains derivatives, so the identity holds as an identity of closed operators with a loss of
radius. For $1<a'<a$, $1<b'<b$, the derivatives are bounded from $Y(a,b)$ to $Y(a',b')$:
\[
  \|\partial_\tau\|\le\frac{q_a}{2(1-q_a)^2},\qquad \|\partial_\xi\|\le\frac{2q_b}{(b'-1)(1-q_b)^2},\qquad
  q_a=a'/a,\ q_b=b'/b,
\]
which for $(a,b)=(65/64,5/4)$, $(a',b')=(129/128,9/8)$ are at most $8385$ and $1440$. Hence $C(s)$ and $M(s)$ are bounded
from the domain of $L$ in $Y(a,b)$ to $Y^\sharp(a',b')$; for $h$ in that domain,
$J^\sharp C(s)h=M(s)L(s)h-(B^\sharp+s)C(s)h\in Y^\sharp(a',b')$, so $C(s)h$ is in the maximal, hence the closed, domain
of $K$ there, and~\eqref{eq:KCML} holds. The inclusion $\ell^1(129/128,9/8)\subset\mathbb T^\sharp(129/128,9/8)$ has norm
$\le1$, since $\kappa_1^m\le\kappa_1^{|m|}$ and $\sqrt{\kappa_2^{2j}+\kappa_2^{-2j}}\le2\kappa_2^j$; the argument $\Dom_{\min}=\Dom_{\max}$
of \cref{app:fredholm} also holds in the torus space. So $C(s)h$ lies in the domain on which the certificates of
\cref{prop:Kinj} assert injectivity.

\subsection{The gauge mode (\cref{eq:gaugemode})}\label{app:gauge}

Let $X_0=e^{\mu\tau}Z$, $Z=\mu^{-1}\partial_\tau+\xi\partial_\xi$, $W=\mu+\xi^2Q$ and $F=\zeta+\log W$. The radial
block of the metric~\eqref{eq:metric} is proportional to $e^{-2\zeta}(u_++u_-)^{-2}du_-du_+$, and $X_0$ translates
$u_\pm$; the Lie derivative gives $\delta\zeta=e^{\mu\tau}(Z\zeta-1)$. The area radius $e^{-\zeta}\xi/W$ is a scalar,
so comparing its variation with its derivative along $X_0$ gives $\delta W=e^{\mu\tau}ZW$, hence
\[
  \delta Q=e^{\mu\tau}(ZQ+2Q),\qquad \delta F=e^{\mu\tau}(ZF-1),\qquad \delta\phi=e^{\mu\tau}Z\phi .
\]
A direct computation gives $[D^\sigma,X_0]=e^{\mu\tau}D^\sigma$. Substituting in~\eqref{eq:omega}, every component
has the same profile: $\delta\omega_i=e^{\mu\tau}(Z+1)\omega_i$ for $i=2,3,4$, and also for $\omega_1=2\xi Q+2\mu\partial_\xi F$,
where the extra terms add up to $\omega_1$. More generally, for every admissible $w$,
\[
  D\Omega(w)\big[e^{\mu\tau}(Z+1)w\big]=e^{\mu\tau}(Z+1)\,\Omega(w),
\]
by the commutator, the quadratic homogeneity of $\Gamma_2$ and the factor $\xi$ in the differential and quadratic
terms. At the exact background the right-hand side vanishes, which gives $L(\mu)g_*=0$ and $C(\mu)g_*=0$.

\section{Explicit formulas}\label{app:formulas}

\subsection{Weights}
The existence proof of~\cite{RT2019} uses $(\kappa_1,\kappa_2)=(65/64,5/4)$. Our strong space is $\Yp=Y(65/64,5/4)$
with component weights $\eta=(916,4096,243,1233)$; the sharp certificates use $(129/128,9/8)$. The real semi-axis
of the Bernstein ellipse of parameter $5/4$ is $\tfrac12(\tfrac54+\tfrac45)=\tfrac{41}{40}$, the radius of the
domain $\Mhat$.

\subsection{The constraint map}
With the encoding $h_i^-=h_i$, $h_i^+=-Ph_i$ ($i=2,3,4$), $D^+=\partial_\tau+\mu(1+\xi)\partial_\xi$, and the
formulas~\eqref{eq:omega} and~\cite[\S2]{RT2019} for $\Omega_1^+$ and $\Omega_{2\sharp}^+$,
\[
  C(s)h=\Big(\tfrac12(1-P)\,\delta\Omega_1^+[h],\ -P\,\delta\Omega_{2\sharp}^+[h]\Big),
\]
where
\begin{align*}
  \delta\Omega_1^+[h]&=\xi(D^++s+\mu)h_1+2\mu h_1+\tfrac14\xi\,\delta\mathcal Q_1[h],\\
  \delta\Omega_{2\sharp}^+[h]&=\xi(D^++s+\mu)h_2^++2\mu(h_2^+-h_3^+)
    +\xi\,\delta\big(-(\omega_2^+)^2+2\omega_3^+\omega_2^+-(\omega_4^+)^2\big)[h],
\end{align*}
$\delta(\cdot)[h]$ denotes the derivative at $\omegabar$ in the direction $h$, and
$\mathcal Q_1=\omega_1^2-2\omega_2^-\omega_1-6\omega_2^+\omega_1+4\omega_3^+\omega_1+(\omega_2^-)^2-2\omega_2^+\omega_2^-
-4\omega_3^+\omega_2^-+(\omega_2^+)^2+4\omega_3^+\omega_2^+-4(\omega_4^+)^2$. In particular $C(s)=C_0+sC_1$ with
\[
  C_1h=\big(0,\ -\xi h_2\big),
\]
since $\xi h_1$ is even and $P\xi P=-\xi$.

\subsection{The operators $K$ and $M$}
With $J_s=\partial_\tau+s+\mu((1+\xi)\partial_\xi+1)$, $X$ multiplication by $\xi$,
$\mathcal G_I=\mathcal G(\omegabar,Ic)$ and $\mathcal G_R=\mathcal G(\omegabar,Rf)$ (\cref{app:KCML}), and using that
$c_1$ and $f_1$ are odd:
\begin{align*}
  (K(s)c)_1&=\big(\partial_\tau+s+\mu(\xi\partial_\xi+1)\big)c_1+(\mathcal G_I)_1,\\
  (K(s)c)_2&=J_sc_2+\mu R_\xi c_1+(\mathcal G_I)_2,\\
  (M(s)f)_1&=-\mu(\xi\partial_\xi+2)Pf_1-(\mathcal G_R)_1,\\
  (M(s)f)_2&=PJ_s(XPf_2)-\mu(Pf_1+Pf_2-f_2+2f_3)-(\mathcal G_R)_2 .
\end{align*}
$\mathcal G_I$ is the bilinear term $\Gamma_2^\sharp(\omegabar,\cdot)$ of~\eqref{eq:sharp}. The first component of $M$ does
not contain $s$, and $M_1f=(0,-\xi f_2)=C_1f$, in agreement with the relation $M_1=C_1$ of \cref{sec:physical-count}.

\section{Table of certificates}\label{app:certificates}

The repository contains, for every step of the proof, the list of files on which the corresponding certificate
depends, with their sizes and sha256 hashes (\texttt{build/reproducao/inventario.json}, rendered in
\texttt{docs/REPRODUCAO\_MAPA.md}). The large files are on Zenodo (\cref{sec:computation-data}), and their hashes are
also in \texttt{build/reproducao/pacote\_dados.json}. The table below gives, for each claim, the check of the
verification program that re-checks it (\cref{sec:computation-verify}), the number of files, and a digest: the
first $16$ hexadecimal digits of the sha256 of the sorted lines ``\texttt{sha256}~~path'' of those files, computed by
\texttt{scripts/digest\_certificados.py}.

\begin{table}[ht]
\centering
\small
% gerado por scripts/digest_certificados.py; nao editar a mao
\begin{tabular}{@{}p{0.33\textwidth}p{0.37\textwidth}rl@{}}
\toprule
claim & level-1 check & files & digest \\
\midrule
data of~\cite{RT2019} (\cref{sec:setting-background}) & download script & 3 & \texttt{75d22969ac3c9d81} \\
\cref{lem:Lreal} & \texttt{L\_real\_toro}, \texttt{radius\_transfer}, \texttt{beta\_plus} & 4 & \texttt{ab394b2cb7a9cca2} \\
\cref{lem:F3}, \cref{thm:counts}~(i) & \texttt{simbolo\_L} & 10 & \texttt{d4a922dbf5a78787} \\
\cref{thm:counts}, A band & \texttt{discos\_A}, \texttt{cobertura\_A}, \texttt{contagem\_A}, \texttt{janelas\_A}, \texttt{deflacao\_q} & 58 & \texttt{c8a9c24901eb3993} \\
\cref{thm:counts}, B band; \cref{prop:nk,prop:witness} at $\sB$ & \texttt{discos\_B}, \texttt{cobertura\_B}, \texttt{contagem\_B}, \texttt{janelas\_B}, \texttt{nk\_raizB} & 86 & \texttt{a62dc9d0d789e2fc} \\
\cref{prop:mu} & \texttt{nk\_s4}, \texttt{identificacao\_s4} & 17 & \texttt{891bdaa6bcbfb3c5} \\
\cref{prop:nk,prop:witness} at $\sK$ & \texttt{nk\_s3b}, \texttt{testemunho\_s3b} & 14 & \texttt{8d44ed8d74426b3a} \\
\cref{prop:Kinj} & \texttt{s3a}, \texttt{rouche\_K}, \texttt{simbolo\_K} & 14 & \texttt{02a66ff6aff02482} \\
\cref{prop:nk} at $\sphys$ & \texttt{nk\_0733} & 19 & \texttt{31458d42757f59dd} \\
\cref{lem:F1} & \texttt{F1} & 1 & \texttt{e4a65f7b6165c1f2} \\
\cref{lem:recovery} & \texttt{R5} & 1 & \texttt{6b2d83307787fd97} \\
\bottomrule
\end{tabular}

\caption{Certificates of the proof, by claim. The digest identifies the set of files on which each claim depends.}
\label{tab:digests}
\end{table}


\bibliographystyle{amsplain}
\bibliography{refs}

\end{document}
